%% LyX 2.4.0-alpha3 created this file.  For more info, see https://www.lyx.org/.
%% Do not edit unless you really know what you are doing.
\documentclass[english,tableposition=top]{report}
\usepackage{lmodern}
\renewcommand{\sfdefault}{lmss}
\renewcommand{\ttdefault}{lmtt}
\usepackage[T1]{fontenc}
\usepackage{textcomp}
\usepackage[latin9]{inputenc}
\setcounter{secnumdepth}{3}
\setcounter{tocdepth}{3}
\usepackage{color}
\definecolor{shadecolor}{rgb}{0.667969, 1, 1}
\usepackage{babel}
\usepackage{array}
\usepackage{float}
\usepackage{booktabs}
\usepackage{framed}
\usepackage{url}
\usepackage{multirow}
\usepackage{amsmath}
\usepackage[unicode=true,pdfusetitle,
 bookmarks=true,bookmarksnumbered=true,bookmarksopen=true,bookmarksopenlevel=2,
 breaklinks=false,pdfborder={0 0 1},backref=section,colorlinks=true]
 {hyperref}

\makeatletter

%%%%%%%%%%%%%%%%%%%%%%%%%%%%%% LyX specific LaTeX commands.
%% Because html converters don't know tabularnewline
\providecommand{\tabularnewline}{\\}

%%%%%%%%%%%%%%%%%%%%%%%%%%%%%% Textclass specific LaTeX commands.
\newenvironment{centred}%
  {\begin{center}\baselineskip=13pt\parskip=1pt}{\end{center}}
\newenvironment{lyxcode}
	{\par\begin{list}{}{
		\setlength{\rightmargin}{\leftmargin}
		\setlength{\listparindent}{0pt}% needed for AMS classes
		\raggedright
		\setlength{\itemsep}{0pt}
		\setlength{\parsep}{0pt}
		\normalfont\ttfamily}%
	 \item[]}
	{\end{list}}

%%%%%%%%%%%%%%%%%%%%%%%%%%%%%% User specified LaTeX commands.
\usepackage{numerica-tables}
\usepackage{numerica-plus}
\newcommand\rel{\,\varrho\;}
\DeclareMathOperator{\erf}{erf}
\DeclareMathOperator{\gd}{gd}
\usepackage{upquote}

\reuse{}

\makeatother

\begin{document}
\title{\texttt{numerica-tables}~\\
}
\author{Andrew Parsloe\\
(\url{ajparsloe@gmail.com})}
\maketitle
\begin{abstract}
The \verb`numerica-tables` package defines a command which enables
the creation of (possibly multi-column) mathematical tables of function
values. \emph{Key=value} settings allow presentation in a wide variety
of table styles within the `formal table' framework of the \texttt{booktabs}
package. \\
\\
\noindent\begin{minipage}[t]{1\columnwidth}%
\begin{shaded}%

\paragraph*{Note:}
\begin{itemize}
\item {\normalsize This document applies to version 2.0.0 of }{\normalsize\texttt{numerica-tables}}{\normalsize .}{\small\par}
\item {\normalsize A version of }{\normalsize\texttt{numerica}}{\normalsize{}
from or later than 2021/12/07 is required; (}{\normalsize\texttt{numerica}}{\normalsize{}
requires }{\normalsize\texttt{amsmath}}{\normalsize , }{\normalsize\texttt{mathtools}}{\normalsize{}
and the \LaTeX 3 bundles }{\normalsize\texttt{l3kernel}}{\normalsize{}
and }{\normalsize\texttt{l3packages}}{\normalsize ).}{\small\par}
\item {\normalsize The }{\normalsize\texttt{booktabs}}{\normalsize{} package
is required.}{\small\par}
\item {\normalsize I refer many times in this document to }{\normalsize\emph{Handbook
of Mathematical Functions}}{\normalsize , edited by Milton Abramowitz
and Irene A. Segun, Dover, 1965. This is abbreviated to }{\normalsize\emph{HMF}}{\normalsize ,
often followed by a reference to a specific table like Table 1.2.}{\small\par}
\item {\normalsize Version 2 of }{\normalsize\texttt{numerica-tables}}{\small\par}
\begin{itemize}
\item {\normalsize is the first stand-alone version (in v.1 of }{\normalsize\texttt{numerica}}{\normalsize{}
the }{\normalsize\texttt{\textbackslash nmcTabulate}}{\normalsize{}
command was available with the }{\normalsize\texttt{tables}}{\normalsize{}
package option);}{\small\par}
\item {\normalsize restricts the third item in }{\normalsize\texttt{rspec}}{\normalsize{}
to }{\normalsize\texttt{rows}}{\normalsize{} only (}{\normalsize\texttt{rstop}}{\normalsize{}
is no longer accepted there);}{\small\par}
\item {\normalsize restricts the third item in }{\normalsize\texttt{cspec}}{\normalsize{}
to }{\normalsize\texttt{cols}}{\normalsize{} only (}{\normalsize\texttt{cstop}}{\normalsize{}
is no longer accepted there);}{\small\par}
\item {\normalsize allows an additional row, if wanted, between table title
and the column header row;}{\small\par}
\item {\normalsize removes the }{\normalsize\texttt{DEL}}{\normalsize{} column
function;}{\small\par}
\item {\normalsize amends documentation.}{\small\par}
\end{itemize}
\end{itemize}
\end{shaded}%
\end{minipage}

\tableofcontents{}
\end{abstract}

\chapter{Introduction}

Entering 
\begin{lyxcode}
~\textbackslash usepackage\{numerica-tables\}
\end{lyxcode}
\noindent in the preamble of a document gives access to a command
\verb`\nmcTabulate` for creating tables of function values in a wide
variety of styles. (Since \verb`numerica-tables` requires \textendash{}
and loads \textendash{} the \verb`numerica` package, all commands
of that package are also available.) All tables are `formal tables'
in the sense of the \texttt{booktabs} package, which is loaded automatically.
Such tables have no vertical rules and few horizontal rules.

\section{Table structure}

I take as my source of models of mathematical tables those presented
in \emph{Handbook of Mathematical Functions}, edited by Milton Abramowitz
and Irene A. Segun, Dover, 1965, not because the typesetting is elegant
(it often isn't) but because \emph{HMF} displays a wide variety of
table styles. The editors of that volume were faced with a host of
different problems requiring a host of different solutions. The \verb`\nmcTabulate`
command aims to reproduce most of those different solutions, within
\verb`booktabs` elegance.

To create a table we need to specify a function to tabulate. The values
this function takes will generally depend on a primary parameter and,
possibly, a number of secondary parameters (which is where much of
the complexity comes from). Mathematical tables are structured in
\emph{columns}. We (nearly always) read \emph{down} a column as the
primary parameter is incremented, generally in regular steps. We need
to decide on the range of values the primary parameter will take and
how fine-grained the tabulation will be \textendash{} what the step
size of its increments will be. Assigning different values to a second
parameter generates a second, third,\ldots{} column. Sometimes rather
than a second parameter, a second, third, \ldots{} function of the
first parameter is tabulated in the successive columns. 

In this document the first parameter is called the \emph{row variable}
\textendash{} its value determines which row we are in; the second
parameter, if present, is called the \emph{column variable \textendash{}
}its value determines which column we are in. A table generally (but
not always) presents the values of the row variable in the first column,
the \emph{row-variable column}, sometimes in distinctive type (e.g.
bolded). The values of the column variable are presented in a \emph{header
row} above the table body of function values. Above the header row
there may be a \emph{title row} and perhaps a \emph{subtitle row}
where other explanatory material can be displayed. Sometimes there
is a \emph{footer row} beneath the table body. Vertical rules are
absent, horizontal rules used sparingly \textendash{} for example,
at the top and bottom of the table, or under the header row, but not
in the body of the table.

\section{Shared syntax}

The \verb`\nmcTabulate` command (short-name form \verb`\tabulate`)
shares the syntax of \verb`\nmcEvaluate` (see \texttt{numerica.pdf}).
When all options are used the command looks like 
\begin{lyxcode}
\noindent \textbackslash nmcTabulate{*}{[}settings{]}\{expr.\}{[}vv-list{]}{[}num.~format{]}
\end{lyxcode}
\begin{enumerate}
\item \verb`*` optional switch; if present ensures a single number output
with no formatting, or an appropriate error message if the single
number cannot be produced; see �\ref{sec:Star-option};
\item \verb`[settings]` comma-separated list of \emph{key=value }settings,
at the heart of creating a table of function values; see Chapter~\ref{chap:=00005CnmcTabulate-settings};
\item \verb`{expr.}` mandatory argument specifying the mathematical expression
or expressions in \LaTeX{} form to be tabulated;
\item \verb`[vv-list]` comma-separated list of \emph{variable=value }items,
in particular containing the initial value of the row variable (and
column variable if one is used); 
\item \verb`[num. format]` optional format specification for presentation
of the numerical results (rounding, padding with zeros, scientific
notation); see �\ref{subsec:Trailing-optional-argument}.
\end{enumerate}
Unlike \verb`\nmcEvaluate` (from\texttt{ }{\ttfamily\verb`numerica`}\texttt{)},
for \verb`\nmcTabulate` 
\begin{itemize}
\item math delimiters are irrelevant \textendash{} it makes no difference
to the display of the result whether the command wraps around math
delimiters, is wrapped within math delimters, or if there are no math
delimiters involved whatever;
\item the two apparently optional arguments straddling the main argument
(\verb`settings` and \verb`vv-list`) are \emph{essential}. Although
both are delimited by square brackets, that is in order to draw on
the code from {\ttfamily\verb`numerica`} for \verb`\nmcEvaluate`.
Each argument contains items \emph{necessary} for the construction
of any table of function values. 
\end{itemize}

\subsection{Settings inherited from \texttt{numerica}}

\label{subsec:Inherited-settings}

Most of the settings available to the command \verb`\nmcEvaluate`
from \verb`numerica` are also available to \verb`\nmcTabulate`.
To save switching between documents I reproduce the table of relevant
options found in \verb`numerica.pdf` (only the punctuation \verb`p`
setting is missing), although for discussion of the options you will
need to refer to that document. (Note that the setting \verb`vvmode`
of v.1 of \verb`numerica` is still available; it is equivalent to
the \verb`vv@` setting.) The \verb`dbg` (and \verb`view`) keys
have been disabled for \verb`\nmcTabulate` at present (but might
be enabled in the future).

\begin{table}[h]
\noindent \centering{}\caption{\protect\label{tab:introSettingsInherited}Settings options inherited
from \texttt{\textbackslash nmcEvaluate}}
\noindent \begin{center}
\begin{tabular}{ll>{\raggedright}p{4cm}>{\raggedright}p{4cm}}
\toprule 
{\small key} & {\small type} & {\small meaning} & {\small default}\tabularnewline
\midrule
{\small\texttt{dbg}} & {\small int} & {\small debug `magic' integer} & {\small\texttt{0}}\tabularnewline
{\small\texttt{view}} &  & {\small equivalent to }{\small\texttt{dbg=1}} & \tabularnewline
{\small\texttt{\textasciicircum}} & {\small char} & {\small exponent mark for sci. notation input} & {\small\texttt{e}}\tabularnewline
{\small\texttt{xx}} & {\small int (}{\small\texttt{0}}{\small /}{\small\texttt{1}}{\small )} & {\small multi-token variable switch} & {\small\texttt{1}}\tabularnewline
{\small\texttt{()}} & {\small int (}{\small\texttt{0}}{\small /}{\small\texttt{1}}{\small /}{\small\texttt{2}}{\small )} & {\small trig. function arg. parsing} & {\small\texttt{0}}\tabularnewline
{\small\texttt{o}} &  & {\small degree switch for trig. functions} & \tabularnewline
{\small\texttt{log}} & {\small num} & {\small base of logarithms for }{\small{\small\verb`\log`}} & {\small\texttt{10}}\tabularnewline
{\small\texttt{vv@}} & {\small int (0/1)} & {\small vv-list calculation mode} & {\small\texttt{0}}\tabularnewline
{\small\texttt{vvd}} & {\small tokens} & {\small vv-list display-style spec.} & {\small\texttt{\{,\}\textbackslash mskip 12mu plus 6mu minus 9mu(vv)}}\tabularnewline
{\small\texttt{vvi}} & {\small tokens} & {\small vv-list text-style spec.} & {\small\texttt{\{,\}\textbackslash mskip 36mu minus 24mu(vv)}}\tabularnewline
{*} &  & {\small suppress equation numbering if }{\small\texttt{\textbackslash\textbackslash}}{\small{}
in }{\small\texttt{vvd}} & \tabularnewline
{\small\texttt{S+}} & {\small int} & {\small extra rounding for stopping criterion for sums} & {\small\texttt{2}}\tabularnewline
{\small\texttt{S?}} & {\small$\text{int}\ge0$} & {\small stopping criterion query terms for sums} & {\small\texttt{0}}\tabularnewline
{\small\texttt{P+}} & {\small int} & {\small extra rounding for stopping criterion for products} & {\small\texttt{2}}\tabularnewline
{\small\texttt{P?}} & {\small$\text{int}\ge0$} & {\small stopping criterion query terms for products} & {\small\texttt{0}}\tabularnewline
\bottomrule
\end{tabular}
\par\end{center}
\end{table}


\chapter{\texttt{\textbackslash nmcTabulate} settings}

\label{chap:=00005CnmcTabulate-settings}In addition to the shared
settings, \verb`\nmcTabulate` has many settings specific to it. They
are discussed in groups in subsequent sections, some in more than
one place. For the main discussion of row-variable settings, see �\ref{sec:Row-variable-settings};
for column-variable settings see �\ref{sec:Column-variable-settings};
for whole-of-table formatting see �\ref{sec:Whole-of-table-formatting};
for formatting the function values in table cells see �\ref{sec:Function-value-formatting}.

\section{Row-variable settings}

\label{sec:Row-variable-settings}
\begin{table}[b]
\centering{}\caption{Row-variable specification}
\begin{center}
\begin{tabular}{ll>{\raggedright}p{4cm}>{\raggedright}p{3cm}}
\toprule 
{\small key} & {\small type} & {\small meaning} & {\small comment}\tabularnewline
\midrule
{\small\texttt{rvar}} & {\small token(s)} & {\small row-variable} & \tabularnewline
{\small\texttt{rstep}} & {\small real num.} & {\small step size} & \tabularnewline
{\small\texttt{rstop}} & {\small real num.} & {\small stop value} & \multirow{2}{3cm}{use only one of {\small\texttt{rstop}}{\small{} or }{\small\texttt{rows}}}\tabularnewline
{\small\texttt{rows}} & {\small int} & {\small number of rows} & \tabularnewline
{\small\texttt{rspec}} & {\small comma list} & {\small\texttt{\{rvar}}{\small , }{\small\texttt{step}}{\small , }{\small\texttt{rows\}}} & {\small short form spec.}\tabularnewline
\bottomrule
\end{tabular}
\par\end{center}
\end{table}
Deciding on a function to tabulate (entered in the main or mandatory
argument of \verb`\nmcTabulate`) will inevitably also mean deciding
on the tabulation variable, the \emph{row} variable, \verb`rvar`,
what value to start tabulating from (which is specified in the vv-list),
what value to tabulate to, \verb`rstop`, and how fine-grained the
tabulation is to be, the step size \verb`rstep`. 

The two tables in the first example below tabulate $\sin x$ and $\cos x$
between $0$ and $1$ in increments of $0.2$. Note that the start
value of the tabulation variable is entered in the vv-list. The reason
for placing it there is that for more complicated functions other
parameters in the function and therefore in the vv-list may depend
on the row variable. Although it will often be the first entry in
the vv-list, it does not need to be. The initial value of the row
variable may depend on other quantities which must necessarily precede
it \textendash{} lie to the right of it \textendash{} in the list.

In the vv-list, the start value of the row variable may be a \LaTeX{}
expression. Both \verb`rstep` and \verb`rstop` can also be \LaTeX{}
expressions. However, they are evaluated \emph{after} the vv-list
is evaluated and so may depend on the values of variables in the vv-list,
including the initial value of the row variable.

The difference in appearance of the two tables below results from
padding with zeros in the second (the asterisk in the trailing optional
argument has the same effect in \verb`\nmcTabulate` as in \verb`\nmcEvaluate`).
As you can see, padding applies not only to the values of the function
but also to the values of the row variable \textendash{} and makes
an obvious improvement to the table's appearance.
\begin{verbatim}
    \tabulate[rvar=x,rstep=0.2,rstop=1]
      { \sin x }[x=0]\qquad
    \tabulate[rvar=x,rstep=0.2,rstop=1]
      { \cos x }[x=0][*]
\end{verbatim}
$\Longrightarrow$     \tabulate[rvar=x,rstep=0.2,rstop=1]
      { \sin x }[x=0] \qquad
    \tabulate[rvar=x,rstep=0.2,rstop=1]
      { \cos x }[x=0][*]\medskip{}

Sometimes (perhaps often) it may prove more convenient to specify
the number of rows, \verb`rows`, explicitly rather than a stop value.
Only one of \verb`rows` and \verb`rstop` should be given, but if
both (inadvertently) are present, it is the value of \verb`rows`
that prevails. The first of the following three tables shows an example
of use.

The second and third tables use an abbreviated form of the row-variable
specification, \verb`rspec`. This is a three-element comma list,
\verb`{rvar,rstep,rows}`. The second table gives a straightforward
example of use. In the third table a \LaTeX{} expression has been inserted
for \verb`rows` in the \verb`rspec` comma list. Like \verb`rstep`
and \verb`rstop`, \verb`rows` can be a \LaTeX{} expression, but it
is evaluated \emph{before} the vv-list and therefore, unlike \verb`rstep`
and \verb`rstop`, cannot depend on quantites specified there like
the row variable.
\begin{verbatim}
    \tabulate[rvar=x,rstep=0.2,rows=6]
      { \sin x/\cos x }[x=0][*] \qquad
    \tabulate[rspec={x,0.2,6}]
      { \tan x }[x=0][*] \qquad
    \tabulate[rspec={x,0.2,1+(1/0.2)}]
      { \sqrt{\sec^2 x - 1} }[x=0][*]
\end{verbatim}
$\Longrightarrow$     \tabulate[rvar=x,rstep=0.2,rows=6]
      { \sin x/\cos x }[x=0][*] \qquad
    \tabulate[rspec={x,0.2,6}]
      { \tan x }[x=0][*] \qquad
    \tabulate[rspec={x,0.2,1+(1/0.2)}]
      { \sqrt{\sec^2 x - 1} }[x=0][*]\medskip{}

\noindent\begin{minipage}[t]{1\columnwidth}%
\begin{shaded}%
In version 1 of \verb`numerica` (when \verb`\nmcTabulate` became
available with the package option \verb`tables`), the third item
in \verb`rspec` was not restricted to \verb`rows` but could also
be \verb`rstop`; parentheses placed around the third item assigned
it to \verb`rows` rather than \verb`rstop`. On reflection, this
now seems obscure and open to error. In v.2 of \verb`numerica-tables`
(in fact its first version as a separate package) the third item of
\verb`rspec` is restricted to \verb`rows` only, which no longer
needs to be parenthesized.\end{shaded}%
\end{minipage}

\subsection{Row-variable column formatting}

\label{subsec:Row-var-col-formatting}The padding option ({*}) of
the trailing optional argument is one way of formatting the row-variable
column, but to how many decimal places? Aligned left or right or centred?
Under what heading \textendash{} the example tables so far have simply
used the row variable for the header? And should the row variable
column be at the left of the table, or the right \textendash{} or
both? These and related questions are answered by assigning values
to the keys listed in Table~\ref{tab:Formatting-the-row-variable}.
\begin{table}[H]
\centering{}\caption{\protect\label{tab:Formatting-the-row-variable}Formatting the row-variable
column}
\begin{center}
\begin{tabular}{ll>{\raggedright}p{4cm}l}
\toprule 
{\small key} & {\small type} & {\small meaning} & {\small default}\tabularnewline
\midrule
{\small\texttt{rround}} & {\small int} & {\small rounding} & {\small\texttt{1}}\tabularnewline
{\small\texttt{ralign}} & {\small char (}{\small\texttt{r/c/l}}{\small )} & {\small horizontal alignment} & {\small\texttt{r}}\tabularnewline
{\small\texttt{rfont}} & {\small chars} & {\small font (}{\small\verb`\math<chars>`}{\small )} & \tabularnewline
{\small\texttt{rhead}} & {\small tokens} & {\small header} & {\small\texttt{rvar}}\tabularnewline
{\small\texttt{rhnudge}} & int & {\small nudge header }{\small{\small\verb`<int>`}}{\small{} mu} & {\small\texttt{0}}\tabularnewline
{\small\texttt{rpos}} & {\small int (}{\small\texttt{0}}{\small\ldots}{\small\texttt{4}}{\small )} & {\small column position(s) } & {\small\texttt{1}}\tabularnewline
{\small\texttt{rvar'}} & {\small tokens} & {\small 2nd row-variable col. spec. } & {\small\texttt{rvar}}\tabularnewline
{\small\texttt{rhead'}} & {\small tokens} & {\small header of 2nd rv col. (if it exists)} & {\small\texttt{rvar'}}\tabularnewline
{\small\texttt{rhnudge'}} & int & {\small nudge 2nd rv col. header }{\small{\small\verb`<int>`}}{\small{}
mu} & {\small\texttt{0}}\tabularnewline
\bottomrule
\end{tabular}
\par\end{center}
\end{table}


\subsubsection{Rounding: \texttt{rround}}

After studying the previous tables, we might decide to adjust the
step size, say from $0.2$ to $0.25$. But changing \texttt{rstep}
to the new value gives a disconcerting result (the first table below).
\texttt{numerica-tables} uses a default rounding value of $1$ for
the row variable and has rounded $0.25$ down to $0.2$ and $0.75$
up to $0.8$ accordingly. The second table corrects matters by adjusting
the row-variable rounding (\texttt{rround}) to \texttt{2}.
\begin{verbatim}
    \tabulate[rvar=x,rstep=0.25,rstop=1]
      { \sin x }[x=0][*] (Eh???) \quad
    \tabulate[rvar=x,rstep=0.25,rstop=1,rround=2]
      { \sin x }[x=0][*]
\end{verbatim}
$\Longrightarrow$     \tabulate[rvar=x,rstep=0.25,rstop=1]
      { \sin x }[x=0][*](Eh???) \quad
    \tabulate[rvar=x,rstep=0.25,rstop=1,rround=2]
      { \sin x }[x=0][*]

\subsubsection{Alignment: \texttt{ralign}}

By default, the alignment of all columns is to the right, as in the
previous examples. This lends itself to neat output when padding with
zeros is activated (the \verb`*` in the trailing argument) and when
some values are negative \textendash{}  minus signs can interfere
with neat output in left or centred alignments. But in a case like
the second table in the last example, you might prefer to centre the
headers for both the row and function-value columns. These alignments
are independently set. For the row-variable column the default alignment
is to the right \texttt{ralign=r}; \texttt{ralign=l} (lowercase L)
aligns entries in the row-variable column to the left, and \texttt{ralign=c}
centres entries in the row-variable column. The tables of the next
example use a \texttt{c} alignment to centre the row-variable column
header. The third of those tables shows how minus signs spoil the
effect.

\subsubsection{Font: \texttt{rfont}}

In the second table below bolding by means of the setting \verb`rfont=bf`
has been applied to emphasize the distinction between the row-variable
values and the function values. Possible values for this key are those
characters that can be adjoined to \verb`\math` to give a meaningful
result. Thus other valid values are \verb`it` (italic), \verb`sf`
(sans serif), \verb`tt` (typewriter); \verb`frak` (Fraktur); also
\verb`rm` (roman) is available, but that is the default.

\subsubsection{Row-variable header: \texttt{rhead}}

In the second and third tables, the header for the row-variable column
has also been bolded. The default header is the row-variable symbol.
That can be replaced by giving a value to the key \texttt{rhead}.
I have used \texttt{rhead=\textbackslash boldsymbol\{x\}} (rather
than \verb`\mathbf{x}`) in order to get an italicized bold symbol.
Note that you do not need to include math delimiters in the specification.
It is assumed that \verb`rhead` will sit between \verb`$ $` delimiters
which are inserted automatically by \verb`numerica-tables`.
\begin{verbatim}
    \tabulate
      [rvar=x,rstep=0.25,rstop=1,
        rround=2,ralign=c]
      { \sin x }[x=0][*]\qquad
    \tabulate
      [rvar=x,rstep=0.25,rstop=1,rround=2,
        ralign=c,rfont=bf,rhead=\boldsymbol{x}]
      { \sin x }[x=0][*]\qquad
    \tabulate
      [rvar=x,rstep=0.25,rstop=0.5,rround=2,
        ralign=c,rfont=bf,rhead=\boldsymbol{x}]
      { \sin x }[x=-0.5][*]
\end{verbatim}
$\Longrightarrow$     \tabulate
      [rvar=x,rstep=0.25,rstop=1,
        rround=2,ralign=c]
      { \sin x }[x=0][*]\qquad
    \tabulate
      [rvar=x,rstep=0.25,rstop=1,
        rround=2,ralign=c,rfont=bf,rhead=\boldsymbol{x}]
      { \sin x }[x=0][*]\qquad
    \tabulate
      [rvar=x,rstep=0.25,rstop=0.5,rround=2,
        ralign=c,rfont=bf,rhead=\boldsymbol{x}]
      { \sin x }[x=-0.5][*]\medskip{}

In these tables the row-variable column has been given a centred alignment.
The third table shows what goes wrong when \emph{some} values in the
row-variable column are negative. Better then is to use padding, a
right alignment (the default), and to use a phantom in the header.
The first table below does this. The second table incorporates kerning
into the header to achieve the same effect:
\begin{verbatim}
    \tabulate
      [rvar=x,rstep=0.25,rstop=0.5,rround=2,
        rfont=bf,rhead=\boldsymbol{x}\hphantom{0}]
      { \sin x }[x=-0.5][*]\qquad
    \tabulate
      [rvar=x,rstep=0.25,rstop=0.5,rround=2,
        rfont=bf,rhead=\boldsymbol{x}\mkern 9 mu]
      { \sin x }[x=-0.5][*]
\end{verbatim}
$\Longrightarrow$     \tabulate
      [rvar=x,rstep=0.25,rstop=0.5,rround=2,
        rfont=bf,rhead=\boldsymbol{x}\hphantom{0}]
      { \sin x }[x=-0.5][*]\qquad
    \tabulate
      [rvar=x,rstep=0.25,rstop=0.5,rround=2,
        rfont=bf,rhead=\boldsymbol{x}\mkern 9 mu]
      { \sin x }[x=-0.5][*]\medskip{}

(To my eye, aligning the $\boldsymbol{x}$ above the first column
of digits after the decimal point gives a better result than truly
centring it in the column; compare these examples with the first two
tables of the previous example.)

\subsubsection{Nudging~the~header: \texttt{rhnudge}}

However, you might prefer to avoid inserting positioning commands
into the actual row-variable header, obscuring its true content. You
can avoid doing this by setting the key \texttt{rhnudge}. 

The first table below reverts to the default right alignment, avoids
any positioning commands in the row-variable header, but instead nudges
it into position with the setting \texttt{rhnudge=9}. For positive
nudge values, nudging works in the \emph{opposite} sense to the alignment.
The units for nudging are mu (math units, 18 to a quad), but only
a number \textendash{} generally an integer \textendash{} should be
specified; the `mu' is supplied by \verb`numerica-tables`.

In the second table below the row variable takes single digit integer
values, while the row-variable name now occupies more than one character.
With a right alignment the header would protrude out to the left.
By giving \texttt{rhnudge} a \emph{negative} value (\texttt{rhnudge=-12}
in the example) it is brought back to a centred position in the row-variable
column. 
\begin{verbatim}
    \tabulate
      [rvar=x,rstep=0.25,rstop=0.5,rround=2,
        rfont=bf,rhead=\boldsymbol{x},rhnudge=9]
      { \sin x }[x=-0.5][4*]\qquad
    \tabulate
      [rvar=x_{\text{int}},rstep=1,rstop=4,
        rround=0,rfont=bf,rhnudge=-12,
        rhead=\boldsymbol{x_{\text{int}}}]
      { \sin x_{\text{int}} }[x_{\text{int}}=0][4*]
\end{verbatim}
$\Longrightarrow$     \tabulate
      [rvar=x,rstep=0.25,rstop=0.5,rround=2,
        rfont=bf,rhead=\boldsymbol{x},rhnudge=9]
      { \sin x }[x=-0.5][4*]\qquad
    \tabulate
      [rvar=x_{\text{int}},rstep=1,rstop=4,
        rround=0,rfont=bf,rhnudge=-12,
        rhead=\boldsymbol{x_{\text{int}}}]
      { \sin x_{\text{int}} }[x_{\text{int}}=0][4*]

\subsubsection{Position in the table: \texttt{rpos}}

\label{subsec:Row-var-col-pos}By default, the row-variable column
is the \emph{first} column of the table. Its position is determined
by the value of the key \texttt{rpos}: 
\begin{itemize}
\item \texttt{rpos=0}, suppressed (no row-variable column); 
\item \texttt{rpos=1}, first column (the default); 
\item \texttt{rpos=2}, last column; 
\item \texttt{rpos=3}, first and last columns;\texttt{ }
\item \texttt{rpos=4}, first and last columns, with the values in the last
column a user-defined function of the first; see �\ref{subsec:Second-row-var-col};
\item Any other integer acts like \texttt{rpos=1}.
\end{itemize}
An example with \texttt{rpos=3} is given shortly below, �\ref{subsec:Multiple-function-tables}. 

\subsubsection{\texttt{rvar'}, \texttt{rhead'}, \texttt{rhnudge'}}

These settings become relevant only when \texttt{rpos=4}; see �\ref{subsec:Second-row-var-col}.

\subsection{Adjoined multi-function tables}

\label{subsec:Adjoining-tables}How might one tabulate multiple functions
simultaneously? \emph{HMF} has many, many examples where multiple
functions (like the trigonometric or the hyperbolic functions) are
tabulated in separate columns of the same table. 

With the settings described so far, one way is to adjoin single column
tables. In the tables below, which display as a single multi-columned
table, I have used three different \texttt{rpos} settings (\texttt{rpos=1}
is implicit in the first). This is one way to build a table that displays
as multi-column. If you use this method, note that the \texttt{\%}
comment characters are essential at the end of the last argument of
the \verb`\tabulate` commands if you want the tables to abut exactly.
Omitting them results in a space between the tables.
\begin{verbatim}
    \tabulate
      [rspec={x,0.2,6}]
      { \sin x }[x=0][*]%
    \tabulate
      [rpos=0,rspec={x,0.2,6}]
      { \cos x }[x=0][*]%
    \tabulate
      [rpos=2,rspec={x,0.2,6}]
      { \tan x }[x=0][*]
\end{verbatim}
$\Longrightarrow$     \tabulate
      [rspec={x,0.2,6}]
      { \sin x }[x=0][*]%
    \tabulate
      [rpos=0,rspec={x,0.2,6}]
      { \cos x }[x=0][*]%
    \tabulate
      [rpos=2,rspec={x,0.2,6}]
      { \tan x }[x=0][*]\medskip{}

However, tabulating more than one function at a time is too common
a need to have to resort to a fudge like adjoining tables. \verb`numerica-tables`
offers a systematic way of doing this; see �\ref{subsec:Multiple-function-tables}.

\section{Column-variable settings}

\label{sec:Column-variable-settings}When a function of \emph{two}
variables is being tabulated, we generally think of one variable as
the primary variable and the other as a parameter. To tabulate such
a function, one way to proceed, is to create and adjoin separate tables,
one per parameter value, but that is clumsy. A more systematic procedure
is to specify, in addition to the row variable, a \emph{column} variable
and its start, step and stop values. 

\begin{table}[t]
\caption{Column-variable specification}

\centering{}\begin{center}
\begin{tabular}{ll>{\raggedright}p{4cm}l}
\toprule 
{\small key} & {\small type} & {\small meaning} & {\small default}\tabularnewline
\midrule
{\small\texttt{cvar}} & {\small token(s)} & {\small column-variable} & \tabularnewline
{\small\texttt{cstep}} & {\small real num.} & {\small step size} & \tabularnewline
{\small\texttt{cstop}} & {\small real num.} & {\small stop value} & {\small either }{\small\texttt{cstop}}\tabularnewline
{\small\texttt{cols}} & {\small int} & {\small number of columns} & {\small or }{\small\texttt{cols}}\tabularnewline
{\small\texttt{cspec}} & {\small comma list} & {\small\texttt{\{cvar,cstep}}{\small ,}{\small\texttt{cols\}}} & {\small short form spec.}\tabularnewline
\bottomrule
\end{tabular}
\par\end{center}
\end{table}
In the following example \verb`cvar=k` is the column variable. I
have chosen a step size \verb`cstep=2` and a stop value \verb`cstop=9`.
As with the row variable, the start value ({\ttfamily\verb`k=3`})
of the column variable is specified in the vv-list. Although in the
example these values are numbers, all three values could be \LaTeX{}
expressions that evaluate to numbers. In particular, the expressions
for step and stop values may include the row and column variables
(in the example $x$ and $k$) which are assigned their initial vv-list
values. Note also the setting for \verb`rhead` which shows the reader
of the table that the numerical values displayed in the column headers
are values of \verb`k`. This usage occurs throughout \emph{HMF}.
\begin{verbatim}
    \tabulate
      [rspec={x,0.2,6},rhead=x\backslash k,
         cvar=k,cstep=2,cstop=9]
      { \sin kx }[k=3,x=0][*]
\end{verbatim}
$\Longrightarrow$     \tabulate
      [rspec={x,0.2,6},rhead=x\backslash k,
         cvar=k,cstep=2,cstop=9]
      { \sin kx }[k=3,x=0][*] \medskip{}

Again, as with the row variable, rather than using an explicit stop
value\texttt{ }{\ttfamily\verb`cstop`}, you might prefer
to specify the number of columns, \verb`cols`, explicitly. I could
have replaced {\ttfamily\verb`cstop=9`} with \verb`cols=4`
to get the same result. Note that the number of columns specified
here is the number of \emph{function-value} columns; the row-variable
column is ignored for this count.

It is worth pointing out explicitly that if \verb`cols` is specified,
then it is possible to have a \emph{zero} step size, \verb`cstep=0`.
An example where this is useful is presented in �\ref{subsec:Footer-functions}.
(A similar comment applies to \verb`rows` and \verb`rstep`.)

And again, as with the row variable, it is possible to condense the
specification into a comma list with the key \verb`cspec`. This is
a $3$-element comma list of the form \verb`{cvar,cstep,cols}`.\footnote{This is a change from v.1 of \texttt{numerica}; see the boxed comment
at the end of �\ref{sec:Row-variable-settings}.} Thus, for the preceding table I could have written
\begin{verbatim}
    \tabulate
      [rspec={x,0.2,6},rhead=x\backslash k,
         cvar=k,cstep=2,cols=4]
      { \sin kx }[k=3,x=0][*]
\end{verbatim}
or
\begin{verbatim}
    \tabulate
      [rspec={x,0.2,6},rhead=x\backslash k,
         cspec={k,2,4}]
      { \sin kx }[k=3,x=0][*]
\end{verbatim}
and produced the same table. 

\verb`cstep`, \verb`cstop` and \verb`cols` can all be \LaTeX{} expressions.
The first two are evaluated \emph{after }the vv-list; \verb`cols`
is evaluated \emph{before }the vv-list. Hence \verb`cstep` and \verb`cstop`
may depend on the row and column variables, which are given their
initial values in the vv-list.

\subsection{Column header formatting}

\label{subsec:Column-header-formatting}There are four built-in style
settings for the header to the column-variable (or function-value)
columns (the `ch' prefix evoking `column header'). If these don't
meet your needs or otherwise satisfy, then it is possible to define
your own header to the function value columns using the key \verb`chead`.
First I discuss the built-in styles.

\subsubsection{Header style: single-column case}

\begin{table}
\centering{}\caption{Formatting the column-variable header}
\begin{center}
\begin{tabular}{ll>{\raggedright}p{4cm}l}
\toprule 
{\small key} & {\small type} & {\small meaning} & {\small default}\tabularnewline
\midrule
{\small\texttt{chstyle}} & {\small int (0\ldots 4)} & {\small header style} & {\small\texttt{0}}\tabularnewline
{\small\texttt{ctitle}} & {\small tokens} & {\small single col. alternative header} & \tabularnewline
{\small\texttt{chead}} & {\small tokens} & {\small user-defined header} & \tabularnewline
{\small\texttt{calign}} & {\small char (r/c/l)} & {\small column alignment} & {\small\texttt{r}}\tabularnewline
{\small\texttt{chnudge}} & {\small int} & {\small nudge header }{\small{\small\verb`int`}}{\small{} mu} & {\small\texttt{0}}\tabularnewline
{\small\texttt{chround}} & {\small int} & {\small column header rounding} & {\small\texttt{0}}\tabularnewline
\bottomrule
\end{tabular}
\par\end{center}
\end{table}
When there is only one column of function values, the function being
tabulated is by default set as the header to the column. This corresponds
to setting \texttt{ctitle={*}} (see �\ref{subsec:Title:-ctitle-setting}
below). You may want some other header. Then give \texttt{ctitle}
some other value (although note that giving it the value \texttt{{*}{*}}
will set both the function and the vv-list as the header; again see
�\ref{subsec:Title:-ctitle-setting}). Whatever value you set, it
will be typeset between math delimiters (\verb`$` signs) and can
be nudged (see �\ref{subsec:Nudgingtheheaders:-chnudge}) left or
right to fine-tune its position in the column. (If you want an asterisk
as the header, you will need to place it between \emph{two} pairs
of braces, \texttt{ctitle=\{\{{*}\}\}}, to prevent it being misinterpreted
as the default setting.)

If you want some more complicated header, perhaps not constrained
by the \verb`$` delimiters, then give \texttt{chead} a value. This
key I discuss below in �\ref{subsec:chead}. \texttt{chead} is entirely
up to the user to specify, including any math environment and positioning. 

If both \texttt{ctitle} and \texttt{chead} are given, the \texttt{chead}
value prevails.

\subsubsection{Header~style: multi-column case}

\texttt{chstyle=0} which is the default gives a header of the form
displayed in the last example, with only the column-variable \emph{value}
at the head of each column. This style generally requires the row-variable
header to indicate what the values denote, as in the example above
where \texttt{rhead=x\textbackslash backslash k}, the backslash separating
row from column variable. \emph{HMF} contains a multitude of instances
of this style; see Tables~9.7, 17.5, 21.1, 24.3, 27.4, etc. for examples.

\texttt{chstyle=1} changes the header of the \emph{first} function
value column to the form \emph{variable=value} \textendash{} in the
example below, to $k=3$. This may be an appropriate choice when a
small rounding value is being used and the resulting columns are narrow.
I can find only one real instance in \emph{HMF}, Table~26.7. Note
that the row-variable setting \texttt{rhead} no longer needs the \texttt{\textbackslash backslash
k} part since the column variable is now explicitly indicated. (The
first table in the example below.)

\texttt{chstyle=2} changes the header of all function-value columns
to the form \emph{variable=value}.\emph{ }In \emph{HMF} examples are
Tables~7.4, 7.9, 10.10, 16.6, etc. Again, the row-variable setting
\texttt{rhead} no longer needs the \texttt{\textbackslash backslash
k} part since the column variable is now explicitly indicated (the
second table in the example).
\begin{verbatim}
    \tabulate
      [rspec={x,0.2,6},
         cspec={k,2,3},chstyle=1]
      { \sin kx }[k=3,x=0][3*]\quad
    \tabulate
      [rspec={x,0.2,6},
         cspec={k,2,3},chstyle=2]
      { \sin kx }[k=3,x=0][3*]
\end{verbatim}
$\Longrightarrow$     \tabulate
      [rspec={x,0.2,6},
         cspec={k,2,3},chstyle=1]
      { \sin kx }[k=3,x=0][3*]\quad
    \tabulate
      [rspec={x,0.2,6},
         cspec={k,2,3},chstyle=2]
      { \sin kx }[k=3,x=0][3*]\medskip{}

\texttt{chstyle=3} fills each column-variable header with the expression
being tabulated but with the column variable replaced by its respective
values. See \emph{HMF} Tables~5.4, 8.1, 9.1, 19.1, etc. for examples.
Note that if the column-variable value is \verb`1`, the \verb`1`
will be displayed:
\begin{verbatim}
    \tabulate
      [rspec={x,0.2,6},
         cspec={k,2,3},chstyle=3]
      { \sin kx }[k=1,x=0][4*]
\end{verbatim}
$\Longrightarrow$     \tabulate
      [rspec={x,0.2,6},
         cspec={k,2,3},chstyle=3]
      { \sin kx }[k=1,x=0][4*]\medskip{}

In this last example you may not want the \verb`1` displayed. To
achieve that result put \texttt{chstyle=4}. This results in a header
as for \texttt{chstyle=3} but if the column variable takes the value
\verb`1`, it has an empty replacement:
\begin{verbatim}
    \tabulate
      [rspec={x,0.2,6},
         cspec={k,2,3},chstyle=4]
      { \sin kx }[k=1,x=0][4*]
\end{verbatim}
$\Longrightarrow$     \tabulate
      [rspec={x,0.2,6},
         cspec={k,2,3},chstyle=4]
      { \sin kx }[k=1,x=0][4*]

\subsubsection{User-defined header: \texttt{chead}}

\label{subsec:chead}If the function in the last example were, for
instance, $k+\sin kx$, then neither replacing $k$ by \verb`1` nor
an empty replacement would be appropriate. In that case the only recourse
is to use the \verb`chead` key. Users can assign whatever value they
like to \verb`chead`. The assignment must contain the correct number
of tab characters (\verb`&`) for the \emph{column-variable columns
only}. It is a header only to the function-value columns. The user
will need to insert \verb`$` signs or other math delimiters as appropriate.
This differs from the practice for \verb`rhead`, but \verb`chead`
is potentially far more complicated. Thus for $k+\sin kx$,
\begin{verbatim}
    \tabulate
      [rspec={x,0.2,6},
         cspec={k,2,3},
         chead=$1+\sin x$&$3+\sin3x$&$5+\sin 5x$]
      { k+\sin kx }[k=1,x=0][4*]
\end{verbatim}
$\Longrightarrow$     \tabulate
      [rspec={x,0.2,6},
         cspec={k,2,3},
         chead=$1+\sin x$&$3+\sin3x$&$5+\sin 5x$]
      { k+\sin kx }[k=1,x=0][4*]\medskip{}

\noindent Non-empty content for the \verb`chead` key overrides any
\verb`chstyle` setting and, in the case of a table with only a single
function-value column, overrides any \verb`ctitle` setting.$ $

\subsubsection{Alignment: \texttt{calign}}

The function-value columns are aligned right, \verb`calign=r`, by
default. Also available are \verb`calign=c` for centred alignment
and \verb`calign=l` (lowercase L) for left alignment. Using centred
alignment with {\ttfamily\verb`chstyle=2`} in a previous
example table gives
\begin{verbatim}
    \tabulate
      [rspec={x,0.2,6},ralign=c,
         cspec={k,2,3},chstyle=2,calign=c]
      { \sin kx }[k=3,x=0][*]
\end{verbatim}
$\Longrightarrow$     \tabulate
      [rspec={x,0.2,6},rhead=x,ralign=c,
         cspec={k,2,3},chstyle=2,calign=c]
      { \sin kx }[k=3,x=0][*]\medskip{}

\noindent The first column of function values looks better, but the
minus signs spoil the effect in the others. Handling signs in tables
is discussed below; see �\ref{subsec:Signs}.

\subsubsection{Nudging~the~headers: \texttt{chnudge}}

\label{subsec:Nudgingtheheaders:-chnudge}In left or right alignment
it is possible to nudge the headers in the opposite direction by giving
a numerical value to the the key \verb`chnudge`. The header is moved
by the specified number of mu (math units; 18 to a quad). Note that
the `mu' does not need to be written. {\ttfamily\verb`numerica-tables`}
provides that. In the next example I have chosen \verb`chnudge=12`
to nudge the column headers to the left to give a centred effect to
the header but leaving the function values with their (potentially)
awkward minus signs right aligned.
\begin{verbatim}
    \tabulate
      [rspec={x,0.2,6},ralign=c,
         cspec={k,2,3},chstyle=2,chnudge=12]
      { \sin kx }[k=3,x=0][*]
\end{verbatim}
$\Longrightarrow$     \tabulate
      [rspec={x,0.2,6},ralign=c,
         cspec={k,2,3},chstyle=2,chnudge=12]
      { \sin kx }[k=3,x=0][*]\medskip{}

The \verb`chnudge` value does not need to be positive. Negative nudges
can be useful when a column header is \emph{longer} than the rounded
function values. In the second example below, I've reduced the rounding
value for function values to $3$, and chosen an initial $k$ value
of $100$ to ensure this circumstance. To centre the column headers
I have used \verb`chnudge=-9`.
\begin{verbatim}
    \tabulate
      [rspec={x,0.2,6},ralign=c,
         cspec={k,2,3},chstyle=2,chnudge=-9]
      { \sin kx }[k=100,x=0][3*]
\end{verbatim}
$\Longrightarrow$     \tabulate
      [rspec={x,0.2,6},ralign=c,
         cspec={k,2,3},chstyle=2,chnudge=-9]
      { \sin kx }[k=100,x=0][3*]

\subsubsection{Rounding: \texttt{chround}}

In the examples so far, the column variable has incremented in integer
steps. The default rounding value for the column variable is $0$
(for the row variable it is $1$), so if it increments by some non-integer
amount, the result will be confusing \textendash{} if $k$ incremented
by, say, $0.25$, starting from $k=3$, then the next column would
also have a header $k=3$ (since $3.25$ with a rounding value $0$
rounds to $3$). The appropriate key to remedy this state of affairs
is {\ttfamily\verb`chround`}. For a step size of $0.25$
the appropriate setting is {\ttfamily\verb`chround=2`}.
\begin{verbatim}
    \tabulate
      [rspec={x,0.2,6},ralign=c,
         cspec={k,0.25,3},chstyle=2,chround=2]
      { \sin kx }[k=3,x=0][*]
\end{verbatim}
$\Longrightarrow$     \tabulate
      [rspec={x,0.2,6},rhead=x,ralign=c,
         cspec={k,0.25,3},chstyle=2,chround=2]
      { \sin kx }[k=3,x=0][*]

\section{Multiple functions in a single table}

\label{subsec:Multiple-function-tables}As already noted in �\ref{subsec:Adjoining-tables},
tabulating more than one function at a time is too common a need to
have to resort to a fudge like adjoining tables. There is a systematic
way of handling this task available in \verb`numerica-tables`. In
v.1 of \verb`numerica`, it sufficed to enter the functions in the
main argument separated by commas, and to \emph{precede the first
function with a comma}, which was the signal \texttt{numerica} needed
to make the internal adjustments for a multi-function table. In v.2
of \verb`numerica-tables` this option is still available, but rather
than use a `trick' like preceding the first function with a comma,
the preferred option now is to use a new (with v.2) setting, \verb`multifn`. 

The first table below uses the old `trick' of starting the main
argument with a comma; the second table uses the \verb`multifn` setting
(and note in both the \verb`o` setting indicating that the arguments
of $\sin$ and $\cos$ are in degrees):
\begin{verbatim}
    \tabulate[o, rround=0,
              rvar=\theta,rstep=10,rstop=90]
      { ,\sin \theta,\cos \theta }[\theta=0][*]
    \quad
    \tabulate[o,multifn,rpos=2,rround=0,
              rvar=\theta,rstep=10,rstop=90]
      { \sin \theta,\cos \theta }[\theta=0][*]
\end{verbatim}
$\Longrightarrow$     \tabulate[o,rvar=\theta,rstep=10,rstop=90,
        rround=0]
  	{ ,\sin \theta,\cos \theta }[\theta=0][*]
    \quad
    \tabulate[o,multifn,rpos=2,rvar=\theta,rstep=10,
        rstop=90,rround=0]
  	{ \sin \theta,\cos \theta }[\theta=0][*]\medskip{}

These tables suggest a space saving possibility: since $\sin$ and
$\cos$ are complementary functions ($\cos\theta=\sin(90-\theta)$),
the values in the bottom half of the table duplicate values in the
top half, only with the columns reversed. This is the reason for the
space saving \texttt{rpos=4} setting (�\ref{subsec:Second-row-var-col})
which enables complementary functions to be tabulated in `half tables'
(for examples see \emph{HMF} Tables 4.10\textendash 4.12 for the trigonometric
functions).

A comma may not always be a convenient separator \textendash{} it
may occur in one of the functions being tabulated (perhaps in \verb`\max`
or \verb`\min`). By assigning a value to the setting \verb`multifn`,
\begin{lyxcode}
multifn=<char>
\end{lyxcode}
the assisgned character can be used to separate the functions. In
the following example, a semicolon is used. Further, the row-variable
column is duplicated on the right by using the \verb`rpos=3` setting,
and the table gives another illustration of the use of \verb`chead`:
\begin{verbatim}
    \tabulate[o,multifn=;,rpos=3,rround=0,
        rvar=\theta,rstep=10,rstop=90,
        chead=max\hphantom{00} & min\hphantom{00}]
      { \max(\sin \theta,\cos \theta);
             \min(\sin \theta,\cos \theta) }
      [\theta=0][*]
\end{verbatim}
$\Longrightarrow$     \tabulate[o,multifn=;,rpos=3,rvar=\theta,
              rstep=10,rstop=90,rround=0,
              chead=max\hphantom{00} & min\hphantom{00}]
  	{ \max(\sin \theta,\cos \theta);
             \min(\sin \theta,\cos \theta) }
      [\theta=0][*]   \medskip{}

The glaring omission in the table is any explicit statement of what
the functions are that are being tabulated. Maximum and minimum, yes,
but of what? That is (potentially) remedied with the \verb`ctitle`
setting discussed next; see �\ref{subsec:Title:-ctitle-setting}.

\section{Whole-of-table formatting}

\label{sec:Whole-of-table-formatting}There are a number of settings
pertaining to the appearance of the table as a whole, things like
the position of the row-variable column, division of the function
values into blocks to aid readability, the presence of horizontal
rules or of a collective column title or of a footer row. I discuss
these here.

\begin{table}[H]
\noindent \centering{}\caption{\protect\label{tab:Table-formatting-settings}Table formatting}
\noindent \begin{center}
\begin{tabular}{ll>{\raggedright}p{4cm}l}
\toprule 
{\small key} & {\small type} & {\small meaning} & {\small default}\tabularnewline
\midrule
{\small\texttt{ctitle}} & {\small token(s)} & {\small collective title for function-value columns} & \tabularnewline
{\small\texttt{cmidrow}} & {\small token(s)} & {\small inter-header/title row for} {\small function-value columns} & \tabularnewline
{\small\texttt{rules}} & {\small char(s)} & {\small horizontal rules template} & {\small\texttt{ThB}}\tabularnewline
{\small\texttt{foot}} & {\small token(s)} & {\small content of footer line} & \tabularnewline
{\small\texttt{rpos}} & {\small int (}{\small\texttt{0}}{\small\ldots}{\small\texttt{4}}{\small )} & {\small row-variable column position(s)} & {\small\texttt{1}}\tabularnewline
{\small\texttt{rbloc}} & {\small comma list} & {\small division of rows into blocks} & \tabularnewline
{\small\texttt{rblocsep}} & {\small length} & {\small extra spacing between blocks of rows} & {\small\texttt{1 ex}}\tabularnewline
\bottomrule
\end{tabular}
\par\end{center}
\end{table}


\subsection{Title for function-value columns: \texttt{ctitle}}

\label{subsec:Title:-ctitle-setting}The function-value columns have
individual headers, formatted in the various ways provided by the
settings discussed in previous sections, but it can also be helpful
to have a collective title for these columns. We saw the need in the
last example. The need is met with the \verb`ctitle` key. This can
be set to whatever you like (e.g. \texttt{ctitle=\textbackslash text\{Fred\}})
but, to more purpose, I use the setting to clarify the last example:
\begin{verbatim}
     \tabulate[o,multifn=;,rpos=3,rround=0,
        rvar=\theta,rstep=10,rstop=90,
        ctitle=\sin\theta{,}\,\cos\theta,
        chead=max\hphantom{00} & min\hphantom{00}]
      { \max(\sin \theta,\cos \theta);
             \min(\sin \theta,\cos \theta) }
      [\theta=0][*]
\end{verbatim}
$\Longrightarrow$      \tabulate[o,multifn=;,rpos=3,rround=0,
        rvar=\theta,rstep=10,rstop=90,
        ctitle=\sin\theta{,}\,\cos\theta,
        chead=max\hphantom{00} & min\hphantom{00}]
      { \max(\sin \theta,\cos \theta);
             \min(\sin \theta,\cos \theta) }
      [\theta=0][*]\medskip{}

\noindent Now it is clearer what is being tabulated, although the
reader is still being asked to interpret rather than read what the
table is showing. 

There are two in-built settings for \verb`ctitle`: \verb`ctitle=*`,
which makes the formula the title, and \verb`ctitle=**`, which makes
a title of the formula and vv-list. Surely, \verb`ctitle=*` is what
we want: 
\begin{verbatim}
    \tabulate[o,multifn=;,rpos=3,rround=0,
        rvar=\theta,rstep=10,rstop=90,
        ctitle=*,
        chead=max\hphantom{00} & min\hphantom{00}]
      { \max(\sin \theta,\cos \theta);
             \min(\sin \theta,\cos \theta) }
      [\theta=0][*]
\end{verbatim}
$\Longrightarrow$     \tabulate[o,multifn=;,rpos=3,rround=0,
        rvar=\theta,rstep=10,rstop=90,
        ctitle=*,
        chead=max\hphantom{00} & min\hphantom{00}]
      { \max(\sin \theta,\cos \theta);
             \min(\sin \theta,\cos \theta) }
      [\theta=0][*]\medskip{}

\noindent Well, that is clear but the length of the title distorts
the table. This is where the \verb`cmidrow` key can help (see below)
but we can also use an \verb`aligned` environment within \verb`ctitle`:
\begin{verbatim}
     \tabulate[o,multifn=;,rpos=3,rround=0,
        rvar=\theta,rstep=10,rstop=90,
        ctitle={\begin{aligned}
                \max(\sin\theta,\cos\theta)\\[-0.7ex]
                \min(\sin\theta,\cos\theta)
               \end{aligned}},
        chead=max\hphantom{00} & min\hphantom{00}]
      { \max(\sin \theta,\cos \theta);
             \min(\sin \theta,\cos \theta) }
      [\theta=0][*]
\end{verbatim}
$\Longrightarrow$      \tabulate[o,multifn=;,rpos=3,rround=0,
        rvar=\theta,rstep=10,rstop=90,
        ctitle={\begin{aligned}
                \max(\sin\theta,\cos\theta)\\[-0.7ex]
                \min(\sin\theta,\cos\theta)
               \end{aligned}},
        chead=max\hphantom{00} & min\hphantom{00}]
      { \max(\sin \theta,\cos \theta);
             \min(\sin \theta,\cos \theta) }
      [\theta=0][*]\medskip{}

The table is no longer distorted in width. Note the \verb`\\[-0.7ex]`
within the \verb`aligned` environment. This shrinks the vertical
space between the two lines of the title. Without it, the lines are
too far apart.

\subsection{Inter-header/title row: \texttt{cmidrow}}

Some tables need to fit more header material or title material into
their rows than can be comfortably accommodated there. For examples,
see \emph{HMF} Tables 7.9 (error function for complex arguments),
17.7 (Jacobian zeta function), 21.1 (eigenvalues of spheroidal wave
functions), and 26.7 (probability integrals). One way of handling
this problem is to resort to more complicated environments in header
and title rows. Another, more direct way, is to insert a row between
the header row and title row. 

I have chosen \verb`cmidrow` for the key name, in the sense of a
row `mid header and title rows'. The initial `c' emphasizes that
it is constrained to the span of the column-variable (or function-value)
columns only (like \verb`chead` and \verb`ctitle`). The entire content
is the responsibility of the user, 
\begin{lyxcode}
cmidrow=<tokens>
\end{lyxcode}
including insertion of the necessary number of tab characters, \verb`&`,
and any math delimiters required. 
\begin{verbatim}
    \tabulate[o,multifn=;,rpos=3,rround=0,
        rvar=\theta,rstep=10,rstop=90,
        ctitle=\max(\sin\theta{,}\cos\theta),
        cmidrow=\multicolumn{2}{c}
                 {$\min(\sin\theta,\cos\theta)$},
        chead=max\hphantom{00} & min\hphantom{00}]
      { \max(\sin \theta,\cos \theta);
            \min(\sin \theta,\cos \theta) }
      [\theta=0][*]
\end{verbatim}
$\Longrightarrow$     \tabulate[o,multifn=;,rpos=3,rround=0,
        rvar=\theta,rstep=10,rstop=90,
        ctitle=\max(\sin\theta{,}\cos\theta),
        cmidrow=\multicolumn{2}{c}
                  {$\min(\sin\theta,\cos\theta)$},
        chead=max\hphantom{00} & min\hphantom{00}]
      { \max(\sin \theta,\cos \theta);
            \min(\sin \theta,\cos \theta) }
      [\theta=0][*]\medskip{}

\noindent The display looks the same as in the previous example but
was obtained perhaps more straightforwardly.

\subsection{Horizontal rules: \texttt{rules}}

\label{subsec:Rules:-rules-setting}The \texttt{booktabs} package
which \texttt{numerica} uses is most emphatic that one should `1.
Never, ever use vertical rules. 2. Never use double rules.' Most
of the tables proper in \emph{HMF} lack rules of any kind although
closer inspection shows smaller tables within the text generally \emph{are}
delimited by horizontal rules (often also with vertical rules).\footnote{The tables in \emph{HMF} are often inelegantly typeset, and sometimes
ugly. For all that, I have used it as a valuable source for the variety
of structures that the editors found necessary, or at least useful,
for presenting a multitude of different kinds of numerical data.} I have used horizontal rules in the various examples in the present
document because these too are tables within text. Some form of delineation
seems necessary. 

\begin{table}
\noindent \centering{}\caption{\protect\label{tab:Rules}Rules. {\small (In the \textquoteleft span\textquoteright{}
column, \textquoteleft f-v\textquoteright =function-value; \textquoteleft r-v\textquoteright =row-variable;
\textquoteleft$<$~table\textquoteright{} indicates that the rule
spans the table but is trimmed at each end.)}}
\noindent \begin{center}
\begin{tabular}{lll>{\raggedright}p{3cm}l}
\toprule 
{\small char} & {\small rule} & {\small position} & {\small span} & {\small default rule thickness}\tabularnewline
\midrule
{\small\texttt{T}} & {\small top} & {\small above table} & {\small table} & {\small\texttt{\textbackslash heavyrulewidth=.08em}}\tabularnewline
{\small\texttt{t}} & {\small title} & {\small below title} & {\small f-v cols} & {\small\texttt{\textbackslash cmidrulewidth =.03em}}\tabularnewline
{\small m} & {\small midrow} & {\small below midrow} & {\small f-v cols (if $1$ r-v col.)}{\small\par}

{\small$<$ table (if $2$ r-v cols)} & {\small\texttt{\textbackslash cmidrulewidth =.03em}}\tabularnewline
{\small\texttt{h}} & {\small header} & {\small below header} & {\small table} & {\small\texttt{\textbackslash lightrulewidth=.05em}}\tabularnewline
{\small\texttt{f}} & {\small footer} & {\small above footer} & {\small table} & {\small\texttt{\textbackslash cmidrulewidth =.03em}}\tabularnewline
{\small\texttt{B}} & {\small bottom} & {\small below table} & {\small table} & {\small\texttt{\textbackslash heavyrulewidth=.08em}}\tabularnewline
\bottomrule
\end{tabular}
\par\end{center}
\end{table}
The \verb`rules` key enables precisely which rules are used to be
specified. The value of the key is a `word' \textendash{} a sequence
of letters \textendash{} where the characters have the significance
and default thicknesses (from \texttt{booktabs}) shown in Table~\ref{tab:Rules}.
The default setting is \texttt{rules=ThB}. To insert a rule beneath
the title, for example, change this to \texttt{rules=TthB}. If in
addition you are using a footer row and want a rule above it, then
the specification is \texttt{rules=TthfB} and if you are using a row
between header and title rows (a `midrow') and want a rule beneath
that too, then the spec. is \texttt{rules=TtmhfB}. To my eye that
is too many rules; at most only one of title and midrow rules should
be used.

The midrow rule changes its behaviour depending on whether there are
two row-variable columns \textendash{} on the left and right of the
table \textendash{} or not. If there is only one row-variable column
then, like the title rule, the midrow rule spans only the function-value
columns. If there are two row-variable columns then the midrow rule
spans the table but is trimmed by 0.5 em at each end. That degree
of trim is the \verb`booktabs` default but can be changed by giving
a different value to \verb`\cmidrulekern` in the preamble, e.g. \verb`\cmidrulekern=1em`.
Note that the changed trim will also apply to the title rule.

If you wish to change the thickness of a rule from its default, then
enter new values for any or all of {\small\texttt{\textbackslash heavyrulewidth}},
{\small\texttt{\textbackslash lightrulewidth}}, {\small\texttt{\textbackslash cmidrulewidth}}
in the preamble. The values listed in Table~\ref{tab:Rules} are
the default values in the \texttt{booktabs} package (except for the
midrow and footer rules, which \texttt{booktabs} does not cover; in
\texttt{numerica-tables} these rules are assigned a thickness of {\small\texttt{\textbackslash cmidrulewidth}}).

In the example table below, a rule for the column title has been specified
(the \verb`t` in the setting \verb`rules=TthB`). Also note the use
of \texttt{ctitle={*}{*}}. The fornula contains an extra parameter
$a$, assigned a value in the vv-list. It now makes sense to display
the vv-list in the column title (but note the braces around \texttt{k}
and \texttt{x} in the vv-list so that they don't display).
\begin{verbatim}
    \tabulate
      [rspec={x,0.25,5},rround=2,rhnudge=9,
         cspec={k,0.25,3},chstyle=2,chround=2,
           ctitle=**,rules=TthB]
      { a\sin kx }[a=2/\pi,{k}=3,{x}=0][*]
\end{verbatim}
$\Longrightarrow$     \tabulate
      [rspec={x,0.25,5},rround=2,rhnudge=9,
         cspec={k,0.25,3},chstyle=2,chround=2,
           ctitle=**,rules=TthB]
      { a\sin kx }[a=2/\pi,{k}=3,{x}=0][*]

\subsection{Footer row: \texttt{foot} setting}

Some tables have a footer row and \verb`numerica-tables` allows such
a row to be inserted, but its entire content, with one exception,
is the responsibility of the user, including insertion of the necessary
number of tab characters \verb`&`. This will be $1$ less than the
total number of columns (including row-variable columns) in the table
\textendash{} or some adjustment thereof if you use \verb`\multicolumn`.
You can put into the footer what you wish: 
\begin{lyxcode}
foot=<tokens>
\end{lyxcode}
(\emph{HMF} uses the footer mainly for cryptic descriptions of the
accuracy and needs of interpolation methods.)

The one exception is when \texttt{foot={*}}. This will fill the footer
with the header, but \emph{reversed}. This is useful for tabulating
complementary functions like the sine and cosine or, more generally,
$f(x)$ and $g(x)$ where $g(x)=f(k-x)$ for some constant $k$. Values
for the complementary function are read from the bottom up and require
a reversed row-variable column on the right of the table; see �\ref{subsec:Second-row-var-col}.

\subsubsection{Footer functions}

\label{subsec:Footer-functions}It is also possible to use the footer
for displaying the values of certain column functions. {\ttfamily\verb`numerica-tables`}
provides four of these. They can be used in the footer (and only in
the footer): \verb`SUM`, \verb`AVE` (average), \verb`MAX` and \verb`MIN`.
These functions act on the function values of the column they are
in. They \emph{do not} combine mathematically: entering \verb`MAX-MIN`
in the footer of a given column will produce a footer entry containing
two values (those of \verb`MAX` and \verb`MIN`) separated by a minus
sign. The numerical output from each function is automatically wrapped
in math delimiters (\verb`$`) so that minus signs display correctly.

\noindent\begin{minipage}[t]{1\columnwidth}%
\begin{shaded}%
In v.1 of \verb`numerica`, tables supported \emph{five} footer functions.
The fifth was \verb`DEL=MAX-MIN`. It is no longer supported. First,
the name was unclear and in any case, it is easy to calculate from
\verb`MAX` and \verb`MIN`.\end{shaded}%
\end{minipage}

In the following example, I have chosen a column variable step size
of zero. This is possible because in the column spec., I have also
specified the exact number of columns. Zeroing the step size means
the same set of figures can be used for the four footer functions
to act on.
\begin{verbatim}
    \tabulate
      [rspec={x,0.25,(5)},rround=2,rhnudge=9,
         cspec={k,0,4},ctitle=**,
         chead=SUM\;&AVE\;&MAX\;&MIN\;,
         rules=TthfB,
         foot=&SUM&AVE&MAX&MIN ]
      { a\sin kx }[a=2/\pi,k=3.5,{x}=0][*]
\end{verbatim}
$\Longrightarrow$     \tabulate
      [rspec={x,0.25,(5)},rround=2,rhnudge=9,
         cspec={k,0,4},calign=r,ctitle=**,
         chead=SUM\;&AVE\;&MAX\;&MIN\;,
         rules=TthfB,
         foot=&SUM&AVE&MAX&MIN ]
      { a\sin kx }[a=2/\pi,k=3.5,{x}=0][*]

\subsection{Second row-variable column: \texttt{rpos=4}}

\label{subsec:Second-row-var-col}In �\ref{subsec:Row-var-col-pos}
I discussed the settings \verb`rpos=0,1,2` and in �\ref{subsec:Multiple-function-tables}
gave an example of using {\ttfamily\verb`rpos=3`} where
repeating the row-variable column on the right is helpful. There is
another value available for this key, \verb`rpos=4`. Like {\ttfamily\verb`rpos=3`}
this adds the row-variable column to both left and right sides of
the table, but for the right column the values are a function of those
in the left column ({\ttfamily\verb`rpos=3`} corresponds
to the function being the identity). The value given to the key \verb`rvar'`
determines the function used and the value given to the key \verb`rhead'`
determines the header for the right-hand row-variable column. If \verb`rhead'`
is omitted it defaults to a blank header, unless the \verb`rvar'`
setting is also omitted, when \verb`rpos=4` behaves like \verb`rpos=3`.

For example, the sine and cosine are complementary functions; when
working in degrees, $\cos\theta=\sin(90-\theta)$. We can exploit
this fact to halve the table size needed to tabulate the two functions.
The example also gives an illustration of the use of an  expression
in the third element of \verb`rspec`.
\begin{verbatim}
    \tabulate[o,multifn,rpos=4,
      rspec={\theta,5,1+45/5},rround=0,
        chnudge=14,rvar'=90-\theta,
          rules=ThfB,foot=*]
      { \sin\theta,\cos\theta }[\theta=0][*]
\end{verbatim}
$\Longrightarrow$     \tabulate[o,multifn,rpos=4,
      rspec={\theta,5,1+45/5},rround=0,
        chnudge=14,rvar'=90-\theta,
          rules=ThfB,foot=*]
      { \sin\theta,\cos\theta }[\theta=0][*]\medskip{}

\noindent The values of sines from $0$ to $45$ degrees are read
downwards from the first column of function values, and from 45 to
90 degrees are read upwards from the second column of function values.
For cosines it is downwards from the second column and upwards from
the first column. The reversed footer line indicates the change of
columns to use. In the example note
\begin{itemize}
\item the setting of \verb`rvar'` to a function (\verb`90-\theta`) of
the row variable (\verb`\theta`); 
\item the blank header for the \verb`rvar'` column (since no value was
set for \verb`rhead'`);
\item the footer setting \verb`foot=*` to obtain the reversed header in
the footer;
\item the rule \emph{above} the footer row specified by the \verb`f` added
to the \verb`rules` setting, \verb`rules=ThfB`.
\end{itemize}
Note also the degree setting \verb`o` in the settings option. 

Although there is a significant space saving with tables like this
(see \emph{HMF} Tables 4.10, 4.11, 4.12), they are not `kind to the
reader'. They require a certain concentration to read and in my view
should be avoided unless space is seriously constrained. \emph{HMF}
Tables 6.1 and 6.2 are tables of the gamma function and its relatives
where $y=x-1$ is used in the row-variable column on the right (stemming
from $y!=\Gamma(x-1)$); \emph{HMF }Table 6.5 in effect uses $\langle1/x\rangle$
(the nearest integer to $1/x$) for the row variable on the right.

\subsection{Separating blocks of rows: \texttt{rbloc}}

Readability of long columns of figures can be aided by breaking the
columns into blocks with extra white space between blocks of rows.
This is achieved with the \texttt{rbloc} key:
\begin{lyxcode}
rbloc~=~<comma~list~of~positive~integers>
\end{lyxcode}
specifies how many rows belong to each block. For example, \verb`rbloc={5,5,6}`
breaks the table into blocks of $5$ rows, $5$ rows, then $6$ rows.
If the number of rows in the table is greater than the sum of the
entries in the comma list, then division into blocks continues as
specified by the last entry in the comma list. Thus \texttt{rbloc=5}
(strictly \texttt{rbloc=\{5\}} but the braces can be omitted in this
case since no comma is enclosed) divides a table into blocks of $5$
rows; \texttt{rbloc=\{1,5\}} divides a table into $1$ row followed
by blocks of $5$ rows. A division of this kind may be appropriate
when, say, the row variable runs from $0$ to $1$ in increments of
$0.1$ \textendash{} there are $11$ rows of which the first (when
the row variable is zero) may have distinctive values. 

\noindent\begin{minipage}[t]{1\columnwidth}%
\begin{shaded}%

\subsubsection*{The pull of the nice round number}

However, this is not how \emph{HMF} sets out its tables.\emph{ }The
dominant practice in \emph{HMF} is division into blocks of (generally)
$5$ rows, many of which start with a zero value for the row variable.
Rather than isolate this initial value, they include it in the first
block of $5$, then continue with blocks of $5$ until a single isolated
row is left at the bottom of the page or the table. There seems to
be a psychological need to finish a page or table with the row variable
set to a nice round number. Thus: tabulate from $0$ to $10$ rather
than $0$ to $9$, from $0$ to $1$ rather than $0$ to $0.9$, and
even from $0$ to $30$ or $0$ to $2$ rather than $0$ to $29$
or $0$ to $1.9$. Using blocks of $5$ the consequence is that there
is always an isolated line at the end \textendash{} a kind of punctuation
mark to signal the end of the page or the table.\end{shaded}%
\end{minipage}

\medskip{}
In the next example I have divided the columns into blocks of $5$
rows by means of the setting \texttt{rbloc=5}.
\begin{verbatim}
    \tabulate[o,rspec={\theta,10,1+90/10},
                          rround=0,rbloc=5]
        { ,\sin\theta, \cos\theta }[\theta=0][*]
\end{verbatim}
$\Longrightarrow$     \tabulate[o,rspec={\theta,10,1+90/10},
                          rround=0,rbloc=5]
		{ ,\sin \theta, \cos \theta}[\theta=0][*]

\subsubsection{Adjusting the extra space\texttt{ rblocsep} }

By default \texttt{numerica} sets the extra space between blocks of
rows at \verb`1 ex`. This value can easily by changed with the setting
\texttt{rblocsep=<length>}. The units need to be included in the specification. 

\section{Formatting function values}

\label{sec:Function-value-formatting}
\begin{table}
\noindent \centering{}\caption{\protect\label{tab:Function-value-formatting}Formatting function
values}
\noindent \begin{center}
\begin{tabular}{cc>{\raggedright}p{4cm}c}
\toprule 
{\small key} & {\small type} & {\small meaning} & {\small default}\tabularnewline
\midrule
{\small\texttt{(pad)}} & {\small int} & {\small (t-notation) phantom padding} & \tabularnewline
{\small\texttt{signs}} & {\small int} & {\small sign handling for function-values} & {\small\texttt{0}}\tabularnewline
{\small\texttt{diffs}} & {\small int} & {\small insert differences \& pre-pad with zeros} & {\small\texttt{0}}\tabularnewline
{\small\texttt{Q?}} & {\small tokens} & {\small special cell conditional} & \tabularnewline
{\small\texttt{A!}} & {\small tokens} & {\small special cell formatting} & \tabularnewline
\bottomrule
\end{tabular}
\par\end{center}
\end{table}
In the examples used so far, function values have been limited to
a narrow range, generally $[-1,1]$. What happens when function values
span orders of magnitude?

\subsection{Trailing optional argument}

\label{subsec:Trailing-optional-argument}The primary tool for function-value
formatting is the trailing optional argument of the \verb`\tabulate`
command where the rounding value is specified, padding with zeros
is set or not (generally \emph{set}), and scientific notation is set
or not. Elegant scientific notation, set with an \texttt{x} in the
trailing optional argument, is generally not appropriate for use in
tables; see the first table below. Repeating the \texttt{x} \textendash{}
\verb`xx` \textendash{} in the trailing optional argument (the second
table) so that scientific notation extends to numbers in the range
$[1,10)$ helps, particularly with the \emph{left} alignment chosen
for the function-value column, but the result is wasteful of space
and the repetition of the `$\times10$' is distracting and would
be more so for a larger table. The \texttt{x} specification should
be used in tables, if at all, only for small tables \textendash{}
a few function values at most.
\begin{verbatim}
    \tabulate[rspec={x,1,2*3+1},rround=0]
      { e^x}[x=-5][*x]\qquad
    \tabulate[rspec={x,1,2*3+1},rround=0,calign=l]
      { e^x}[x=-3][*xx]
\end{verbatim}
$\Longrightarrow$     \tabulate[rspec={x,1,2*3+1},rround=0]
      { e^x}[x=-3][*x]\qquad
    \tabulate[rspec={x,1,2*3+1},rround=0,calign=l]
      { e^x}[x=-3][*xx]

\subsubsection{The \texttt{t} option}

\emph{HMF} uses a special notation for coping with function values
spanning orders of magnitude. This notation can be invoked by inserting
\texttt{t} in the trailing optional argument. Repeating the previous
two tables, and adding a \texttt{chnudge} value, gives a more compact
and visually appealing result:
\begin{verbatim}
    \tabulate[rspec={x,1,2*3+1},rround=0,chnudge=24]
      { e^x}[x=-3][*t]\qquad
    \tabulate[rspec={x,1,2*3+1},rround=0,chnudge=24]
      { e^x}[x=-3][*tt]
\end{verbatim}
$\Longrightarrow$     \tabulate[rspec={x,1,2*3+1},rround=0,chnudge=24]
      { e^x}[x=-3][*t]\qquad
    \tabulate[rspec={x,1,2*3+1},rround=0,chnudge=24]
      { e^x}[x=-3][*tt]

\subsection{Padding the exponent: \texttt{(pad)}}

In the second table of the last example some might quibble at the
lack of alignment of the left parentheses. \emph{HMF} tends to align
these and \verb`numerica-tables` offers the setting 
\begin{lyxcode}
(pad)~=~<integer>
\end{lyxcode}
to achieve the effect. (The parentheses are part of the key \textendash{}
a reminder of the \verb`t`-form of scientific notation.) \texttt{<integer>}
is the number of digits/characters to pad to. Repeating the last two
tables with the setting \texttt{(pad)=2} produces the following results:
\begin{verbatim}
    \tabulate[rspec={x,1,2*3+1},rround=0,
                       chnudge=24,(pad)=2]
      { e^x}[x=-3][*t]\qquad
    \tabulate[rspec={x,1,2*3+1},rround=0,
                       chnudge=24,(pad)=2]
      { e^x}[x=-3][*tt]
\end{verbatim}
$\Longrightarrow$     \tabulate[rspec={x,1,2*3+1},rround=0,
                       chnudge=24,(pad)=2]
      { e^x}[x=-3][*t]\qquad
    \tabulate[rspec={x,1,2*3+1},rround=0,
                       chnudge=24,(pad)=2]
      { e^x}[x=-3][*tt]\medskip{}

Note that this setting is relevant only when the \texttt{t} option
is used in the trailing number-formatting argument of the \verb`\tabulate`
command. Examples in \emph{HMF} of the style exemplified by the first
table are, among others, Tables 8.6, 9.2, 20.1, and of the style exemplified
by the second table, among many, Tables 9.9, 10.5, 13.1, 14.1, 19.1.

\subsection{Accommodating signs: \texttt{signs}}

\label{subsec:Signs}Instead of $e^{x}$ as the test function, use
$e^{x}-1.$ Now there are positive, zero and negative function values
to contend with. Recall that in the \verb`t`-notation the \emph{exponent}
is the parenthesized integer part of a number, the \emph{significand}
the following decimal figures. \verb`numerica-tables` offers the
\verb`signs` key to align (or not) the exponents. The setting is
\begin{lyxcode}
signs~=~<integer>
\end{lyxcode}
There are four effective values for \texttt{<integer>} and the do-nothing
default (\texttt{signs=0}):
\begin{itemize}
\item \texttt{signs=2 }inserts a $+$ sign between exponent and significand
of every non-negative number;
\item \texttt{signs=1 }inserts a $+$ sign between exponent and significand
of every non-negative number that immediately precedes or follows
a negative number;
\item \texttt{signs=-1 }inserts a $+$ sign between exponent and significand
of any non-negative number that immediately precedes or follows a
negative number, and inserts a \emph{phantom} $+$ sign between exponent
and significand of every other non-negative number;
\item \texttt{signs=-2 }inserts a \emph{phantom} $+$ sign between exponent
and significand of every non-negative number;
\end{itemize}
In the following examples with \verb`signs=-2`, \verb`signs=-1`
and \verb`signs=2`, all give acceptable results.
\begin{verbatim}
    \tabulate[rspec={x,1,2*3+1},rround=0,
        (pad)=2,signs=-2]
      { e^x-1}[x=-3][4*tt]\qquad
    \tabulate[rspec={x,1,2*3+1},rround=0,
        (pad)=2,signs=-1]
      { e^x-1}[x=-3][4*tt]\qquad
    \tabulate[rspec={x,1,2*3+1},rround=0,
        (pad)=2,signs=2]
      { e^x-1}[x=-3][4*tt]
\end{verbatim}
$\Longrightarrow$     \tabulate[rspec={x,1,2*3+1},rround=0,
        (pad)=2,signs=-2]
      { e^x-1}[x=-3][4*t']\qquad
    \tabulate[rspec={x,1,2*3+1},rround=0,
        (pad)=2,signs=-1]
      { e^x-1}[x=-3][4*t']\qquad
    \tabulate[rspec={x,1,2*3+1},rround=0,
        (pad)=2,signs=2]
      { e^x-1}[x=-3][4*t']\medskip{}

\noindent In \emph{HMF} Table 23.2 illustrates \verb`signs=-2`; Tables
10.1, 13.1, 14.1, 19.1 among many others illustrate \verb`signs=-1`;
and Tables 9.4, 10.6, 20.2, 22.11 among others illustrate \verb`signs=2`.

\verb`signs=1`, however, is an inappropriate setting for these function
values in the \verb`t`-notation:
\begin{verbatim}
    \tabulate[rspec={x,1,2*3+1},rround=0,
        (pad)=2,signs=1]
      { e^x-1}[x=-3][4*tt] \qquad
\end{verbatim}
$\Longrightarrow$     \tabulate[rspec={x,1,2*3+1},rround=0,
        (pad)=2,signs=1]
      { e^x-1}[x=-3][4*tt]\medskip{}

But the \verb`signs` key is not limited to the \verb`t`-notation.
In the following tables where the notation is not used, positive values
for the key, including \verb`signs=1`, give good results (I've included
also the default setting \textendash{} the third table):
\begin{verbatim}
    \tabulate[rspec={x,0.1,9},(pad)=2,signs=2]
      { 10\sin 5x}[x=-0.4][*4]\qquad
    \tabulate[rspec={x,0.1,9},(pad)=2,signs=1]
      { 10\sin 5x}[x=-0.4][*4]\qquad
    \tabulate[rspec={x,0.1,9},(pad)=2]
      { 10\sin 5x}[x=-0.4][*4]
\end{verbatim}
$\Longrightarrow$     \tabulate[rspec={x,0.1,9},(pad)=2,signs=2]
      { 10\sin 5x}[x=-0.4][*4]\qquad
    \tabulate[rspec={x,0.1,9},(pad)=2,signs=1]
      { 10\sin 5x}[x=-0.4][*4]\qquad
    \tabulate[rspec={x,0.1,9},(pad)=2]
      { 10\sin 5x}[x=-0.4][*4]\medskip{}

\noindent \emph{HMF} seems to use \verb`signs=2` when the sign of
the function values changes every few entries and \verb`signs=1`
when there are runs of entries of the same sign. Over the range tabulated
here for $10\sin5x$, they would use the middle table of the three,
\verb`signs=1`. 

\subsection{Differences: \texttt{diffs}}

In fine-grained tables where function values change only slowly from
entry to entry it can be helpful to include a difference entry between
function-value entries as an aid to interpolation (and a test of eyesight).
By entering 
\begin{lyxcode}
diffs~=~<non-negative~integer>
\end{lyxcode}
the \verb`\tabulate` command will include differences in a table.
The \texttt{<non-negative integer>} is the maximum number of digits
in a difference. 
\begin{verbatim}
    \tabulate[rspec={x,0.01,1+(1.05-1)/0.01},rround=2,
      rhnudge=9,chnudge=21,diffs=3]
        { \sinh x }[x=1][*4]
\end{verbatim}
$\Longrightarrow$     \tabulate[rspec={x,0.01,1+(1.05-1)/0.01},rround=2,
      rhnudge=9,chnudge=21,diffs=3]
        { \sinh x }[x=1][*4]\medskip{}

I have deliberately chosen the function and settings here \textendash{}
particularly \texttt{diffs=3} \textendash{} to give a good result.
With the default right alignment of the function-value columns, it
is easy to get this wrong. The evidence will be either in the misalignment
of the first row of function values or unnecessary padding of differences
with leading zeros. It is a good idea to create your table first,
see how function values change between successive rows and judge how
many digits there will be in a difference. In the following examples
I have deliberately put \texttt{diffs=2} and \texttt{diffs=4} to show
the effect of a misjudgement. In the first table the first row of
function values is misaligned by one character. (\texttt{diffs=1}
would have produced a two-character misalignment.) In the second table
the unnecessary fourth digit for the differences results in pre-padding
with $0$. 

In the second table the function $-\sinh x$ is \emph{decreasing},
showing how it is the \emph{absolute value of the difference} between
successive function values that is tabulated. A difference is always
a non-negative value.
\begin{verbatim}
    \tabulate[rspec={x,0.01,1+(1.05-1)/0.01},rround=2,
      rhnudge=9,chnudge=21,diffs=2]
        { \sinh x }[x=1][*4]\qquad
    \tabulate[rspec={x,0.01,1.05},rround=2,
      rhnudge=9,chnudge=30,diffs=4]
        { -\sinh x }[x=1][*4]
\end{verbatim}
$\Longrightarrow$     \tabulate[rspec={x,0.01,1+(1.05-1)/0.01},rround=2,
      rhnudge=9,chnudge=21,diffs=2]
        { \sinh x }[x=1][*4]\qquad
    \tabulate[rspec={x,0.01,1+(1.05-1)/0.01},rround=2,
      rhnudge=9,chnudge=21,diffs=4]
        { -\sinh x }[x=1][*4]\medskip{}

\noindent When the \texttt{diffs} setting is too small, function values
in the first row are misaligned, the amount depending on how much
too small. (A left alignment of the function value column is another
way of tackling this issue.) When the \texttt{diffs} setting is too
big, alignment is fine but differences are padded with unnecessary
leading zeros, meaning the column header will need a bigger nudge
to bring \emph{it} into alignment. 

\subsection{Formatting special values: \texttt{Q?} and \texttt{A!}}

You may wish to highlight or display in some special way a particular
function value or values. \verb`\nmcTabulate` has two related settings
that enable this: \texttt{Q?=<tokens>} and \texttt{A!=<tokens>}. As
the names suggest: Question? and Answer! 

The question should be an expression that \texttt{l3fp} can digest
and produce a boolean answer to (1 for `true' or 0 for `false').
\emph{This is not a \LaTeX{} expression}; this is an \verb`l3fp` expression.\footnote{Documentation about \texttt{l3fp }can be found in \texttt{interface3.pdf},
which is part of the \texttt{l3kernel} bundle.} \verb`numerica-tables` uses \texttt{@} to denote the current function
value, so queries like \texttt{Q?=@<0} (Is the current function value
negative?) or \texttt{Q?=\{}@\texttt{>=pi\}} (Is the current function
value greater than or equal to $\pi$?) are valid questions. (Note
the braces in the second question, used to hide the equality sign.)
Other possible useful components of such questions are \texttt{exp(1)}
for the number $e$, || for logical Or, \texttt{\&\&} for logical
And, and \texttt{!} for logical Not,\texttt{ }as well as the familiar
arithmetic symbols, \texttt{+},\texttt{ -},\texttt{ {*}},\texttt{
/} and \texttt{\textasciicircum}, relation symbols \texttt{<},\texttt{
>},\texttt{ =} and their combinations like \texttt{!=},\texttt{ >=},\texttt{
<=} etc., and parentheses. In addition to these components, \texttt{numerica-tables}
offers \texttt{MAX} and \texttt{MIN} which are the maximum and minimum
function values tabulated, so that, e.g., \texttt{Q?=\{@=MIN\}} (note
the braces) is the question: Is the current function value equal to
the minimum function value for the whole table?

The answer must be in the form of a \LaTeXe{} formatting statement,
again using \texttt{@} to denote the current function value. Thus
\texttt{A!=\textbackslash mathbf\{@\}} is a valid answer; so is \texttt{A!=\textbackslash color\{red\}\{@\}}
(provided you have \texttt{\textbackslash usepackage\{color\}} in
the preamble); and so is \texttt{A!=(@)}. Another valid answer is
\texttt{A!=} , meaning that function values satisfying the \texttt{Q?}
question are omitted from the output. 

This can be useful to suppress `irrelevant' values in a particular
context. For example, suppose we wish to focus on the values of $\cos(m\pi/n)$
lying between $0$ and $\tfrac{1}{2}$ inclusive for certain values
of $m$ and $n$. Rather than cluttering the table with values outside
that interval, we suppress them (the two occurrences of `\texttt{1e-14}'
in the query are there to prevent rounding errors confusing the result):
\begin{verbatim}
    \tabulate
      [rspec={n,1,1+(15-4)},rround=0,rpos=2,rules=Tth,
         cspec={m,1,1+(5-2)},ctitle=*,chstyle=2,
           Q?={@<-1e-14||@>0.5+1e-14},A!=]
      { \cos(m\pi/n) }[n=4,m=2][*4]
\end{verbatim}
$\Longrightarrow$      \tabulate
      [rspec={n,1,1+(15-4)},rround=0,rpos=2,rules=Tth,
         cspec={m,1,1+(5-2)},ctitle=*,chstyle=2,
           Q?={@<-1e-14||@>0.5+1e-14},A!=]
      { \cos(m\pi/n) }[n=4,m=2][*4]

\section{Star option: \texttt{\textbackslash nmcTabulate{*}}}

\label{sec:Star-option}If the \verb`Q?` question is satisfied by
at least one function value then adding a star (asterisk) to the \verb`\tabulate`
command will display the first such instance. Like other starred commands
in the \verb`numerica` suite (\verb`\eval*`, \verb`info*`, \verb`\macros*`,
\verb`\constants*`, \verb`\iter*`, \verb`\solve*` and \verb`\recur*`),
 \verb`\tabulate*` outputs a single number. Using the star means
you do not need an answering \verb`A!` to the query \verb`Q?` since
no formatting of table values is involved.
\begin{verbatim}
    \tabulate*
      [rspec={n,1,1+(15-4)},cspec={m,1,1+(5-2)},
           Q?={@<-1e-14||@>0.5+1e-14}]
      { \cos(m\pi/n) }[n=4,m=2][*4]
\end{verbatim}
$\Longrightarrow$      \tabulate*
      [rspec={n,1,1+(15-4)},cspec={m,1,1+(5-2)},
           Q?={@<-1e-14||@>0.5+1e-14}]
      { \cos(m\pi/n) }[n=4,m=2][*4]. Indeed, if you omit the \verb`Q?` and \verb`A!` settings from
the previous table so that all function values are visible then this
is the value that follows $0.5000$ in the \texttt{m=2} column \textendash{}
the first function value encountered either less than $0$ or greater
than $0.5$.

\subsection{Errors}

If \emph{no} function value satisfies the query then a message is
generated:
\begin{verbatim}
    \tabulate*
      [rspec={n,1,1+(15-4)},cspec={m,1,1+(5-2)},
           Q?=@>1]
      { \cos(m\pi/n) }[n=4,m=2][*4]
\end{verbatim}
$\Longrightarrow$      \tabulate*
      [rspec={n,1,1+(15-4)},cspec={m,1,1+(5-2)},
           Q?=@>1]
      { \cos(m\pi/n) }[n=4,m=2][*4]

And if there is no query at all when the star option is chosen, another
message is shown:
\begin{verbatim}
    \tabulate*
      [rspec={n,1,1+(15-4)},cspec={m,1,1+(5-2)}]
      { \cos(m\pi/n) }[n=4,m=2][*4]
\end{verbatim}
$\Longrightarrow$      \tabulate*
      [rspec={n,1,1+(15-4)},cspec={m,1,1+(5-2)}]
      { \cos(m\pi/n) }[n=4,m=2][*4]

\subsection{Scientific notation}

If you want the number output in scientific notation when the star
option is chosen, then enter the exponent mark in the trailing number-format
option. This is straightforward for a letter like the commonly used
\texttt{e}, but remember that if it is the \texttt{x} option that
you enter, then you will need to place the \verb`\tabulate*` command
between math delimiters, otherwise the \verb`\times` symbol resulting
from the \texttt{x} option will generate a \LaTeX{} error (`Missing
\$ inserted'):
\begin{verbatim}
    $ 
    \tabulate*
      [rspec={n,1,1+(15-4)},cspec={m,1,1+(5-2)},
           Q?={@<-1e-14||@>0.5+1e-14},A!=]
      { \cos(m\pi/n) }[n=4,m=2][*4x]
    $
\end{verbatim}
$\Longrightarrow$     $ 
    \tabulate*
      [rspec={n,1,1+(15-4)},cspec={m,1,1+(5-2)},
           Q?={@<-1e-14||@>0.5+1e-14},A!=]
      { \cos(m\pi/n) }[n=4,m=2][*4x]
    $.

\subsection{Nesting}

A \verb`\tabulate` command can be nested within other commands from
the wider \verb`numerica` package, and those other commands can be
nested within a \verb`\tabulate` command.

\subsubsection{Nesting of \texttt{\textbackslash nmcTabulate}}

Occasionally one might want to extract a value \textendash{} perhaps
a maximum or minimum \textendash{} from a table to insert in another
command. This can be done by nesting a \verb`\tabulate*` command
with an appropriate \verb`Q?` setting within the other command. In
fact, from v.2 of \verb`numerica` we can omit (or forget to include)
the star. All we require is that the \verb`Q?` setting is satisfied
by at least one function value. 
\begin{verbatim}
    \eval{$ 
      (\tabulate
        [rspec={n,1,15},cspec={m,1,5},
             Q?={@=MAX}]
      { \cos(m\pi/n) }[n=4,m=2][*4])\sinh t + 
     (\tabulate
        [rspec={n,1,15},cspec={m,1,5},
             Q?={@=MIN}]
      { \sin(m\pi/n) }[n=4,m=2][*4])\cosh t
         $}[t=2][4]
\end{verbatim}
$\Longrightarrow$     \eval{$ 
      (\tabulate
        [rspec={n,1,15},cspec={m,1,5},
             Q?={@=MAX}]
      { \cos(m\pi/n) }[n=4,m=2][*4])\sinh t + 
     (\tabulate
        [rspec={n,1,15},cspec={m,1,5},
             Q?={@=MIN}]
      { \sin(m\pi/n) }[n=4,m=2][*4])\cosh t
         $}[t=2][4]. 

Evaluating the tables
\begin{verbatim}
     \tabulate
        [rspec={n,1,15},rround=0,rpos=2,rules=Tth,
           cspec={m,1,5},ctitle=*,chstyle=2]
      { \cos(m\pi/n) }[n=4,m=2][*4]
\end{verbatim}
for the cosine and
\begin{verbatim}
    \tabulate
        [rspec={n,1,15},rround=0,rpos=2,rules=Tth,
           cspec={m,1,5},ctitle=*,chstyle=2]
      { \sin(m\pi/n) }[n=4,m=2][*4]
\end{verbatim}
for the sine and checking the entries shows that indeed the maximum
and minimum values are $0.9397$ and $-1.0000$ respectively.

If the \verb`Q?` setting is not satisfied by any function value an
error message is shown; e.g.,
\begin{verbatim}
    \eval{$ 
      (\tabulate
        [rspec={n,1,15},cspec={m,1,5},
          Q?=@>2]
            { \cos(m\pi/n) }[n=4,m=2][*4])\sinh t 
          $}[t=2][4]
\end{verbatim}
$\Longrightarrow$     \eval{$ 
      (\tabulate
        [rspec={n,1,15},cspec={m,1,5},
          Q?=@>2]
            { \cos(m\pi/n) }[n=4,m=2][*4])\sinh t 
          $}[t=2][4]

\noindent Here, `settings (1)' tells us that the message refers
to the (or a) command at the first level of nesting.

\subsubsection{Nesting within \texttt{\textbackslash nmcTabulate}}

Perhaps a more likely situation is to want to nest other commands
within a \verb`\tabulate` command, in particular those from the associated
package \verb`numerica-` \verb`plus`. In the documentation for that
package, as an illustration of the use of those commands, I describe
how they can be used to numerically investigate the timing of signals
between points fixed on a rotating disk. and in particular how to
look for three-point circuits with the property that, despite the
rotation, signals take the same time traversing the circuit in one
direction as the other. 

I show a complicated expression below involving a \verb`\tabulate`
command wrapped around a \verb`\solve` command (from \verb`numerica-plus`)
wrapped around four \verb`\iter*` commands (also from \verb`numerica-plus`).
I also show the table resulting from it all. On my High St laptop
the table takes perhaps two minutes to compile \textendash{} I haven't
timed it exactly but it is well over a minute. Two minutes is the
blink of an eye in a lifetime but an age sitting staring at a computer
screen. Hence I compiled the table separately, saved it to file using
the \verb`\reuse` command (see the discussion below, �\ref{sec:reuse-setting})
and  pasted the saved result into this document.

To explain the expression: suppose the three points on the disk have
polar coordinates $(r,\pm\theta)$ and $(a,0)$ in a co-rotating coordinate
system. In the underlying inertial system, the signal speed is $c$
and the angular velocity $\omega$. The cosine rule for solving triangles
gives for the travel times $t$ between the points 
\[
t=c^{-1}\sqrt{r^{2}+a^{2}-2ra\cos(\theta\pm\omega t)}
\]
 and
\[
t=c^{-1}r\sqrt{2-2\cos(2\theta\pm\omega t)}
\]
where the plus sign is for a signal in the direction of rotation and
the minus sign for a signal against the rotation. Thus the travel
times are fixed points of the expressions on the right. Fixed points
are found with the \verb`\iter*` command, which is where those terms
of the expression come from. The \verb`\solve` command is then used
to find, for a given value of $\theta$, a value of $r$ for which
the difference of the travel times is zero for the circuit traversed
in the two directions. Finally, the \verb`\tabulate` command creates
a table of $r$ values for different values of the row-variable $\theta$
and column variable $a$. 

As noted earlier, it takes a long time to compile this expression
on an `ordinary' laptop, and so I have compiled the expression separately
and pasted the result into this document:
\begin{verbatim}
    
    \tabulate[rspec={\theta,0.2,5},chstyle=2,
              cspec={a,5,3},ctitle=r_{\Delta t=0}]
    { \solve[var=r,vvi=,+=1]
        {% circuit 1231
          2\times\iter*[var=t,+=1]{ 
            c^{-1}\sqrt{a^2+r^2-
              2ar\cos(\theta-\omega t)} }[4]
          + \iter*[var=t,+=1]{ 
              c^{-1}r\sqrt{2-2\cos(2\theta+\omega t)} }[4]
         % circuit 1321
          - 2\times\iter*[var=t,+=1]{ 
              c^{-1}\sqrt{a^2+r^2-
                2ar\cos(\theta+\omega t)} }[4]
          - \iter*[var=t,+=1]{ 
              c^{-1}r\sqrt{2-2\cos(2\theta-\omega t)} }[4]
        }[3]
    }[ c=30,r=a+10,a=10,
            \theta=0.2,\omega=0.2,t=1 ][3*]
\end{verbatim}
$\Longrightarrow$ \begin {tabular}{rccc}
\toprule 
&\multicolumn {3}{c}{$r_{\Delta t=0}$}\tabularnewline 
$\theta $&$a=10$&$a=15$&$a=20$\tabularnewline 
\midrule 
${0.2}$&$10.21$&$15.31$&$20.43$\tabularnewline 
${0.4}$&$10.87$&$16.33$&$21.82$\tabularnewline 
${0.6}$&$12.15$&$18.31$&$24.55$\tabularnewline 
${0.8}$&$14.46$&$21.88$&$29.59$\tabularnewline 
${1.0}$&$18.82$&$28.86$&$39.82$\tabularnewline 
\bottomrule 
\end {tabular}
\begin{verbatim}

\end{verbatim}

\section{Table placement}

\label{sec:Table-placement}There is only one setting in this category
that is part of \texttt{numerica-tables} as such, the tabular vertical
alignment option; see �\ref{subsec:Verticalalignment}. But \LaTeX{}
allows one to insert vertical space with its \verb`\bigskip`, \verb`\medskip`,
\verb`\smallskip`, usually about one line space, a half line space,
and a quarter line space (with stretch and shrink), and \verb`booktabs`
provides \verb`\abovetopsep` and \verb`\belowbottomsep`, both set
by default to \verb`0ex` (or any other unit you care to use) and
easily changed by writing, e.g., \verb`\abovetopsep=1.25ex` if you
want to insert \verb`1.25ex` of space above the table (perhaps to
fit captions).

\subsection{Vertical~alignment}

\label{subsec:Verticalalignment}By writing\texttt{ }
\begin{lyxcode}
valign~=~<char>~
\end{lyxcode}
where \texttt{<char>} is one of \texttt{t}, \texttt{m} or \texttt{b}
the vertical alignment of the table can be set relative to the text
baseline.\texttt{ valign=t }aligns the top of the table with the text
baseline, \texttt{valign=b} the bottom of the table with the text
baseline, \texttt{valign=m} aligns the middle of the table with the
text baseline. By default \texttt{valign=m} is set. Repeating an example
from earlier (�\ref{sec:Row-variable-settings}) I have added letters
A, B, C to show where the baseline is. In the first table the top
of the table aligns with the baseline; in the second table (default
case) the middle of the table aligns with the baseline; in the third
table, the bottom of the table aligns with the baseline.
\begin{verbatim}
    A \tabulate[valign=t,rvar=x,rstep=0.2,rows=6]
      { \sin x/\cos x }[x=0][*] \qquad
    B \tabulate[rspec={x,0.2,1+(1/0.2)}]
      { \tan x }[x=0][*] \qquad
    C \tabulate[valign=b,rspec={x,0.2,(6)}]
      { \sqrt{\sec^2 x - 1} }[x=0][*]
\end{verbatim}
$\Longrightarrow$     A \tabulate[valign=t,rvar=x,rstep=0.2,rows=6,rstop=2]
      { \sin x/\cos x }[x=0][*] \qquad
    B \tabulate[rspec={x,0.2,1+(1/0.2)}]
      { \tan x }[x=0][*] \qquad
    C \tabulate[valign=b,rspec={x,0.2,6}]
      { \sqrt{\sec^2 x - 1} }[x=0][*]\medskip{}

As explained in �\ref{subsec:Adjoining-tables}, tables can be adjoined
to give the appearance of a single larger table. If tables with different
numbers of rows are adjoined in this manner, then a middle alignment
fails and a top alignment is necessary (so that the header rows of
the tables align).

\section{The \texttt{reuse} setting}

\label{sec:reuse-setting}By entering 
\begin{lyxcode}
reuse~=~<non-negative~integer>
\end{lyxcode}
it is possible to specify what is saved when the \verb`\tabulate`
command is followed by a \verb`\reuse` command.
\begin{itemize}
\item \texttt{reuse=0} saves the table as displayed (the default);
\item \texttt{reuse=n} saves the $n$-th \emph{column} of function values
\begin{itemize}
\item if \texttt{rpos} is non-zero: in a comma-separated list of braced
pairs \texttt{\{row-variable value,function value}\};
\item if \texttt{rpos=0}: in a comma-separated list of function values.
\end{itemize}
\item \texttt{reuse=-n} saves the $n$-th \emph{row} of function values
in a comma-separated list that includes the row-variable value in
its appropriate position(s) if \texttt{rpos} is non-zero.
\end{itemize}
In the following example the third row is saved to the control sequence
\verb`\rowiii` by using the setting \verb`reuse=-3`.
\begin{verbatim}
    \tabulate
      [rspec={x,0.25,5},rround=2,rhead=x,
        ralign=r,rhnudge=9,
         cspec={k,0.25,3},chstyle=2,
          chround=2,calign=r,ctitle=**,
           rules=TthB,reuse=-3]
      { a\sin kx }[a=2/\pi,{k}=3,{x}=0][*] 
    \reuse{rowiii}
\end{verbatim}
$\Longrightarrow$     \tabulate
      [rspec={x,0.25,(5)},rround=2,
         rhead=x,ralign=r,rhnudge=9,
         cspec={k,0.25,(3)},chstyle=2,
         chround=2,calign=r,ctitle=**,rules=TthB,reuse=-3]
      { a\sin kx }[a=2/\pi,{k}=3,{x}=0][*] 
    \reuse[renew]{rowiii}\medskip{}

\noindent Now test the content of the control sequence
\begin{lyxcode}
\noindent \verb`\rowiii`~$\Longrightarrow$~\rowiii ~
\end{lyxcode}
\noindent You can see that indeed the third row has been `captured
for posterity'.

Alternatively, we could save the second column of function values
in the control sequence \verb`\colii`:
\begin{verbatim}
    \tabulate[rspec={x,0.25,5},rround=2,rhead=x,
     ralign=r,rhnudge=9,
      cspec={k,0.25,3},chstyle=2,
       chround=2,calign=r,ctitle=**,
        rules=TthB,reuse=2]
         { a\sin kx }[a=2/\pi,{k}=3,{x}=0][*] 
    \reuse{colii}
\end{verbatim}
$\Longrightarrow$     \tabulate
      [rspec={x,0.25,(5)},rround=2,
         rhead=x,ralign=r,rhnudge=9,
         cspec={k,0.25,(3)},chstyle=2,
         chround=2,calign=r,ctitle=**,
         rules=TthB,reuse=2]
      { a\sin kx }[a=2/\pi,{k}=3,{x}=0][*] 
    \reuse[renew]{colii}\medskip{}

Now to see what is saved in \verb`\colii` I use \TeX 's \verb`\meaning`
command to show that it is indeed braced pairs: 
\begin{verbatim}
    \meaning \colii
\end{verbatim}
$\Longrightarrow$ \meaning\colii

Finally, let's use the default \verb`reuse` setting (\verb`reuse=0`)
by not entering an explicit \verb`reuse` setting at all. This should
save the whole table:
\begin{verbatim}
    \tabulate[rspec={x,0.25,5},rround=2,rhead=x,
     ralign=r,rhnudge=9,
      cspec={k,0.25,3},chstyle=2,
       chround=2,calign=r,ctitle=**,
        rules=TthB]
         { a\sin kx }[a=2/\pi,{k}=3,{x}=0][*] 
    \reuse{wholetable}
\end{verbatim}
$\Longrightarrow$ \tabulate[rspec={x,0.25,5},rround=2,rhead=x,
     ralign=r,rhnudge=9,
      cspec={k,0.25,3},chstyle=2,
       chround=2,calign=r,ctitle=**,
        rules=TthB]
         { a\sin kx }[a=2/\pi,{k}=3,{x}=0][*] 
    \reuse[renew]{wholetable}\medskip{}
 

\noindent Now check that the whole table has been saved:\medskip{}

\noindent \verb`\wholetable` $\Longrightarrow$\wholetable\medskip{}

\noindent and, indeed it has.

\chapter{Reference summary}

\section{Commands defined in \texttt{numerica-tables}}

\texttt{\textbackslash nmcTabulate, \textbackslash tabulate}

\section{Settings for \texttt{\textbackslash nmcTabulate}}

\subsubsection*{Row-variable specification �\ref{sec:Row-variable-settings}}
\begin{center}
\begin{center}
\begin{tabular}{ll>{\raggedright}p{4cm}>{\raggedright}p{3cm}}
\toprule 
{\small key} & {\small type} & {\small meaning} & {\small comment}\tabularnewline
\midrule
{\small\texttt{rvar}} & {\small token(s)} & {\small row-variable} & \tabularnewline
{\small\texttt{rstep}} & {\small real num} & {\small step size} & \tabularnewline
{\small\texttt{rstop}} & {\small real num} & {\small stop value} & \multirow{2}{3cm}{either {\small\texttt{rstop}}{\small{} or }{\small\texttt{rows}}{\small ,
not both}}\tabularnewline
{\small\texttt{rows}} & {\small int} & {\small number of rows} & \tabularnewline
{\small\texttt{rspec}} & {\small comma list} & {\small\texttt{\{start}}{\small , }{\small\texttt{step}}{\small , }{\small\texttt{rows\}}} & {\small short form spec.}\tabularnewline
\bottomrule
\end{tabular}
\par\end{center}
\par\end{center}

\subsubsection*{Row-variable column formatting �\ref{subsec:Row-var-col-formatting}}
\begin{center}
\begin{center}
\begin{tabular}{ll>{\raggedright}p{4cm}l}
\toprule 
{\small key} & {\small type} & {\small meaning} & {\small default}\tabularnewline
\midrule
{\small\texttt{rround}} & {\small int} & {\small rounding} & {\small\texttt{1}}\tabularnewline
{\small\texttt{ralign}} & {\small char (}{\small\texttt{r/c/l}}{\small )} & {\small horizontal alignment} & {\small\texttt{r}}\tabularnewline
{\small\texttt{rfont}} & {\small chars} & {\small font (}{\small\verb`\math<chars>`}{\small )} & \tabularnewline
{\small\texttt{rhead}} & {\small tokens} & {\small header} & {\small\texttt{rvar}}\tabularnewline
{\small\texttt{rhnudge}} & int & {\small nudge header }{\small{\small\verb`rhnudge`}}{\small{}
mu} & {\small\texttt{0}}\tabularnewline
{\small\texttt{rpos}} & {\small int (}{\small\texttt{0}}{\small\ldots}{\small\texttt{4}}{\small )} & {\small column position(s)} & {\small\texttt{1}}\tabularnewline
{\small\texttt{rvar'}} & {\small token(s)} & {\small 2nd row-variable col. spec.} & {\small\texttt{rvar}}\tabularnewline
{\small\texttt{rhead'}} & {\small token(s)} & {\small header of 2nd r-v col. (if it exists)} & {\small\texttt{rvar'}}\tabularnewline
{\small\texttt{rhnudge'}} & int & {\small nudge 2nd r-v col. header }{\small{\small\verb`rhnudge'`}}{\small{}
mu} & {\small\texttt{0}}\tabularnewline
\bottomrule
\end{tabular}
\par\end{center}
\par\end{center}

\subsubsection*{Column-variable specification �\ref{sec:Column-variable-settings}.}
\begin{center}
\begin{center}
\begin{tabular}{ll>{\raggedright}p{4cm}l}
\toprule 
{\small key} & {\small type} & {\small meaning} & {\small default}\tabularnewline
\midrule
{\small\texttt{cvar}} & {\small token(s)} & {\small column-variable} & \tabularnewline
{\small\texttt{cstep}} & {\small real num} & {\small step size} & \tabularnewline
{\small\texttt{cstop}} & {\small real num} & {\small stop value} & {\small either }{\small\texttt{cstop}}\tabularnewline
{\small\texttt{cols}} & {\small int} & {\small number of columns} & {\small or }{\small\texttt{cols}}{\small , not both}\tabularnewline
{\small\texttt{cspec}} & {\small comma list} & {\small\texttt{\{cvar,cstep}}{\small ,}{\small\texttt{cols\}}} & {\small short form spec.}\tabularnewline
\bottomrule
\end{tabular}
\par\end{center}
\par\end{center}

\subsubsection*{Column-variable header formatting �\ref{subsec:Column-header-formatting}.}
\begin{center}
\begin{center}
\begin{tabular}{ll>{\raggedright}p{4cm}l}
\toprule 
{\small key} & {\small type} & {\small meaning} & {\small default}\tabularnewline
\midrule
{\small\texttt{chstyle}} & {\small int (}{\small\texttt{0}}{\small\ldots}{\small\texttt{4}}{\small )} & {\small header style} & {\small\texttt{0}}\tabularnewline
{\small\texttt{ctitle}} & {\small token(s)} & {\small single col. alternative header} & \tabularnewline
{\small\texttt{chead}} & {\small token(s)} & {\small user-defined header} & \tabularnewline
{\small\texttt{calign}} & {\small char (r/c/l)} & {\small column alignment} & {\small\texttt{r}}\tabularnewline
{\small\texttt{chnudge}} & {\small int} & {\small nudge header }{\small{\small\verb`chnudge`}}{\small{}
mu} & {\small\texttt{0}}\tabularnewline
{\small\texttt{chround}} & {\small int} & {\small column header rounding} & {\small\texttt{0}}\tabularnewline
\bottomrule
\end{tabular}
\par\end{center}
\par\end{center}

\subsubsection*{Function-value formatting �\ref{sec:Function-value-formatting}.}
\noindent \begin{center}
\noindent \begin{center}
\begin{tabular}{cc>{\raggedright}p{4cm}c}
\toprule 
{\small key} & {\small type} & {\small meaning} & {\small default}\tabularnewline
\midrule
{\small\texttt{(pad)}} & {\small int} & {\small (t-notation) phantom padding} & \tabularnewline
{\small\texttt{signs}} & {\small int} & {\small sign handling for function-values} & {\small\texttt{0}}\tabularnewline
{\small\texttt{diffs}} & {\small int} & {\small insert differences \& pre-pad with zeros} & {\small\texttt{0}}\tabularnewline
{\small\texttt{Q?}} & {\small tokens} & {\small special cell conditional} & \tabularnewline
{\small\texttt{A!}} & {\small token(s)} & {\small special cell formatting} & \tabularnewline
\bottomrule
\end{tabular}
\par\end{center}
\par\end{center}

\subsubsection*{\newpage Whole-of-table formatting �\ref{sec:Whole-of-table-formatting}.}
\noindent \begin{center}
\noindent \begin{center}
\begin{tabular}{ll>{\raggedright}p{4cm}l}
\toprule 
{\small key} & {\small type} & {\small meaning} & {\small default}\tabularnewline
\midrule
{\small\texttt{ctitle}} & {\small token(s)} & {\small collective title for function-value columns} & \tabularnewline
{\small\texttt{cmidrow}}  & {\small token(s) } & {\small inter-header/title row} & \tabularnewline
{\small\texttt{rules}} & {\small char(s)} & {\small horizontal rules template} & {\small\texttt{ThB}}\tabularnewline
{\small\texttt{foot}} & {\small token(s)} & {\small content of footer line} & \tabularnewline
{\small\texttt{rpos}} & {\small int (}{\small\texttt{0}}{\small\ldots}{\small\texttt{4}}{\small )} & {\small row-variable column position(s)} & {\small\texttt{1}}\tabularnewline
{\small\texttt{rbloc}} & {\small comma list} & {\small division of rows into blocks} & \tabularnewline
{\small\texttt{rblocsep}} & {\small length} & {\small extra spacing between blocks of rows} & {\small\texttt{1 ex}}\tabularnewline
\bottomrule
\end{tabular}
\par\end{center}
\par\end{center}

\subsubsection*{Miscellaneous settings �\ref{subsec:Multiple-function-tables}, \ref{subsec:Verticalalignment}.}
\noindent \begin{center}
\noindent \begin{center}
\begin{tabular}{cc>{\raggedright}p{4cm}c}
\toprule 
{\small key} & {\small type} & {\small meaning} & {\small default}\tabularnewline
\midrule
{\small\texttt{multifn}} & {\small char} & {\small multi-function separator} & {\small\texttt{,}}\tabularnewline
{\small\texttt{valign}} & {\small char (t/m/b)} & {\small vertical alignment} & {\small\texttt{m}}\tabularnewline
\bottomrule
\end{tabular}
\par\end{center}
\par\end{center}
\end{document}
